\documentclass[12pt]{article}

% Page layout
\usepackage[margin=1in]{geometry}
\usepackage{setspace}
\onehalfspacing

% Math
\usepackage{amsmath}
\usepackage{amssymb}
\usepackage{amsthm}
\usepackage{mathtools}

% Tables
\usepackage{booktabs}
\usepackage{multirow}
\usepackage{tabularx}
\usepackage{longtable}
\usepackage{array}

% General
\usepackage{graphicx}
\usepackage{hyperref}
\usepackage{fancyvrb}
\usepackage{xcolor}
\usepackage{fancyhdr}
\usepackage{calc}
\usepackage{etoolbox}
\fancyhead{}
\renewcommand{\headrulewidth}{0pt}
\fancyfoot{}

% Theorem environments
\newtheorem{theorem}{Theorem}[section]
\newtheorem{lemma}[theorem]{Lemma}
\newtheorem{proposition}[theorem]{Proposition}
\newtheorem{corollary}[theorem]{Corollary}
\newtheorem{definition}[theorem]{Definition}
\newtheorem*{remark}{Remark}
\newtheorem*{example}{Example}

% Verbatim / code
\DefineVerbatimEnvironment{Highlighting}{Verbatim}{commandchars=\\\{\},fontsize=\small}
\newenvironment{Shaded}{}{}
\newcommand{\NormalTok}[1]{#1}
\newcommand{\KeywordTok}[1]{\textbf{#1}}
\newcommand{\StringTok}[1]{#1}
\newcommand{\CommentTok}[1]{\textit{#1}}
\newcommand{\FunctionTok}[1]{#1}
\newcommand{\OperatorTok}[1]{#1}
\newcommand{\OtherTok}[1]{#1}
\newcommand{\DataTypeTok}[1]{#1}
\newcommand{\BuiltInTok}[1]{#1}

% Misc
\providecommand{\tightlist}{\setlength{\itemsep}{0pt}\setlength{\parskip}{0pt}}
\setlength{\emergencystretch}{3em}
\makeatletter
\patchcmd\longtable{\par}{\if@noskipsec\mbox{}\fi\par}{}{}
\makeatother

\newenvironment{CSLReferences}[2]{}{}
\newcommand{\CSLBlock}[1]{#1}
\newcommand{\CSLLeftMargin}[1]{#1}
\newcommand{\CSLRightInline}[1]{#1}
\newcommand{\CSLIndent}[1]{\hspace{\csllabelwidth}#1}
\begin{document}

\title{Calculus from First Principles: A Limit-Based Symbolic Computation Engine via Log-Puiseux Series}
\author{Jason Parker}
\date{}

\maketitle
\thispagestyle{fancy}

\begin{abstract}
\noindent We describe a symbolic calculus engine in which differentiation,
integration, and limit evaluation are derived from a single underlying
representation: the \textbf{log-Puiseux series}, finite formal sums of
terms \(c \cdot h^p \cdot \log(h)^k\) with \(c \in \bar{\mathbb{Q}}\),
\(p \in \mathbb{Q}\), \(k \in \mathbb{N}\). We establish closure of this
class under addition, multiplication, and composition within the
single-logarithm tower, and show it is strictly more expressive than
pure Puiseux series. Differentiation is implemented as extraction of the
\(h^1\) coefficient from the symbolic expansion of \(f(x+h)\); limits
are read from the leading term, returning certified results of type
\(\mathtt{Exists}\), \(\mathtt{Pole}\), \(\mathtt{DoesNotExist}\), or
\(\mathtt{LimitError}\); integration proceeds via the Risch algorithm
over the differential field the representation naturally provides,
returning \(\mathtt{NonElementary}\) as a proved certificate when no
elementary antiderivative exists. Multivariate calculus follows without
additional machinery. The product rule and chain rule are not
implemented as rewriting rules; both are consequences of the limit
definition and are observable in the output.
\end{abstract}


\section{Introduction}\label{introduction}

Symbolic calculus systems and real analysis share a subject but not a
foundation. In analysis, the derivative is defined as a limit; the
product rule, chain rule, and every rule of differentiation are theorems
proved from that definition. In symbolic computation, the derivative is
defined by those same theorems, recast as rewriting rules. The limit
definition is present in the mathematical background but absent from the
computation. This inversion is not a flaw --- it is a deliberate
foundational choice, made because rewriting rules are efficient and
because the theorems, once proved, carry all the information needed for
the algebraic manipulations that CAS systems perform. But it is a
choice, and it is worth asking what is lost by making it.

What is lost, we argue, is a certain kind of generality. A rule-table
system knows what it has been told. Adding the sine integral
\(\mathrm{Si}(x)\) requires adding rules for its derivative, patterns
for its recognition in integration, and special cases wherever the
system must reason about it. The system grows by declaration. There is
no sense in which \(\mathrm{Si}(x)\) follows from the rules already
present; it must be announced.

This paper describes \textbf{limcalc}, a symbolic calculus engine built
on the opposite principle: the limit definition is the computation, not
the motivation for a computation that proceeds by other means.
Differentiation, integration, and limit evaluation are all consequences
of a single underlying representation --- the \textbf{log-Puiseux
series} --- chosen because it is the natural algebraic object for
expanding \(f(x+h)\) in powers of \(h\). The derivative is the
coefficient of \(h^1\) in that expansion. The limit is the constant
term. We never taught the system the product rule or the chain rule.
Both can be observed using the system, because the system computes
limits the way analysts think about them.

We wish to be precise about what we are and are not claiming. We are not
claiming that limcalc is more capable than existing computer algebra
systems. Mathematica, Maple, and SymPy are mature systems representing
decades of engineering; no research prototype competes with them on
breadth or performance. We are claiming something more specific: that
the limit-first architecture is more mathematically honest. The product
rule is a theorem; our system treats it as one. The chain rule is a
theorem; our system treats it as one. When no elementary antiderivative
exists, the system returns a proved certificate of that fact rather than
a failure. The architecture reflects the logical structure of the
subject as analysts understand it, and the consequences of that
alignment are the subject of this paper.


\subsection{The Classical Approach and Its Foundational
Choice}\label{the-classical-approach-and-its-foundational-choice}

The dominant paradigm in symbolic computation treats calculus as a
branch of algebra. A computer algebra system maintains a symbolic
expression tree and applies rewriting rules to transform it.
Differentiation is defined by a table: \(\frac{d}{dx} x^n = nx^{n-1}\),
\(\frac{d}{dx} \sin(x) = \cos(x)\), \(\frac{d}{dx} e^x = e^x\), together
with the chain rule as a structural rule applied when a composition is
recognized. Integration is handled by a combination of pattern matching,
heuristic transformations, and --- at the deepest level --- the Risch
algorithm operating over differential fields. The historical development
of this tradition, and the engineering tradeoffs that shaped it, are
surveyed by Moses \cite{moses1971}; standard textbook treatments appear
in Davenport, Siret, and Tournier \cite{davenport1988} and Geddes,
Czapor, and Labahn \cite{geddes1992}.

This approach is not wrong. It is a deliberate and well-motivated
foundational choice: represent calculus operations as transformations on
symbolic expressions, and let the algebraic structure of the expression
language carry the mathematical content. The result is systems of
remarkable generality and power. The choice has a cost that is rarely
made explicit: the architecture is inherently extensional. Adding a new
special function --- say, the sine integral \(\mathrm{Si}(x)\) ---
requires adding new rewriting rules for its derivative, new patterns for
its integration, and new cases in every part of the system that must
recognize it. The system grows by accumulation. There is no sense in
which \(\mathrm{Si}(x)\) follows from anything already present; it must
be declared.

The extensional character of rule-table CAS is a practical engineering
reality, not a criticism. But it is worth naming, because it is
precisely what the present work departs from.


\subsection{Automatic
Differentiation}\label{automatic-differentiation}

A different tradition attacks the problem numerically. Automatic
differentiation (AD) --- in both forward mode, via dual numbers, and
reverse mode, via adjoint accumulation --- computes derivatives of
programs to machine precision without symbolic expression manipulation
(Griewank and Walther 2008). Forward-mode AD is, in a sense, closer to
the analyst's view: it propagates the chain rule through actual
computation, maintaining a derivative alongside each value. The result
is exact (within floating-point arithmetic) and efficient.

The foundational choice in AD is to operate on values rather than
expressions. This is the source of both its strength --- no symbolic
explosion, no simplification problem --- and its limitation for our
purposes. AD produces numbers, not symbolic forms. It cannot tell you
that \(\int e^{-x^2}\,dx\) has no elementary closed form; it cannot
return \(\mathrm{erf}(x)\); it cannot state that
\(\lim_{x \to 0} \frac{\sin x}{x} = 1\) as a certified symbolic fact.
The goals are different.


\subsection{Differential Algebra and the Risch
Algorithm}\label{differential-algebra-and-the-risch-algorithm}

The most mathematically sophisticated tradition in symbolic integration
is that of differential algebra, originating with Liouville's theorem on
elementary integrals and formalized computationally by Risch (Risch
1969). The algebraic foundations of differential fields and their
extensions were developed by Kolchin \cite{kolchin1973}. Bronstein's
treatise (Bronstein 2005) remains the definitive computational reference. In this framework, one works in a differential field --- a
field equipped with a derivation satisfying the product rule --- and
integration is the problem of finding an antiderivative within a tower
of field extensions. The Risch algorithm decides, in finite time,
whether an elementary antiderivative exists, and finds it if so.

This tradition is both rigorous and complete for the integration
problem. Our Risch implementation stands directly on Bronstein's work,
and we make no claim to have improved on the algorithm itself. What we
note is that in the differential algebra tradition, as in the rule-table
tradition, the limit definition does not appear as a computational
primitive. The derivation in a differential field is an algebraic object
satisfying axioms; it is not computed as a limit. The connection to the
analytic definition is present in the mathematical background but absent
from the computation.


\subsection{Interval Arithmetic and Rigorous
Numerics}\label{interval-arithmetic-and-rigorous-numerics}

A fourth tradition prioritizes certified numerical results. Systems such
as iRRAM (Müller 2001) and libraries built on MPFR compute with real
numbers to arbitrary precision, with rigorous error bounds. Limits can
be evaluated to any desired precision; derivatives can be approximated
via finite differences with controlled error. The foundational choice
here is to represent real numbers as convergent sequences of rational
approximations rather than as symbolic expressions. The results are
certified but numerical --- a decimal approximation with a proved error
bound, not a symbolic closed form.


\subsection{The Present Work}\label{the-present-work}

limcalc makes a different foundational choice from each of the above.
The starting point is the analyst's definition: the derivative of \(f\)
at \(x\) is

\[f'(x) = \lim_{h \to 0} \frac{f(x+h) - f(x)}{h}\]

and this limit is to be computed symbolically, not approximated
numerically or reduced to a lookup. The key insight is that this
computation is tractable if one chooses the right representation for the
formal expansion of \(f(x+h)\) in powers of \(h\). For a wide class of
functions --- all elementary functions and the standard special
functions --- \(f(x+h)\) can be expanded as a \textbf{log-Puiseux
series} in \(h\): a finite formal sum of terms of the form

\[c \cdot h^p \cdot \log(h)^k\]

where \(c\) is an algebraic number coefficient, \(p \in \mathbb{Q}\) is
the power exponent, and \(k \in \mathbb{N}\) is the logarithmic
exponent. The derivative is then the coefficient of \(h^1\) in this
expansion. The limit \(\lim_{h\to 0} f(x+h)\) is the constant term.
Integration is performed by the Risch algorithm operating on the same
representation.

The log-Puiseux type is the primary novel contribution of this paper.
Pure Puiseux series --- the \(k = 0\) case --- are insufficient: the
cosine integral \(\mathrm{Ci}(x)\) and the exponential integral
\(\mathrm{Ei}(x)\) have genuine \(\log(h)\) singularities near \(x = 0\)
that cannot be represented as pure power series. The log-Puiseux type is
closed under addition and multiplication:

\[\left(c_1 h^{p_1} \log(h)^{k_1}\right) \cdot \left(c_2 h^{p_2} \log(h)^{k_2}\right) = c_1 c_2 \, h^{p_1 + p_2} \log(h)^{k_1 + k_2}\]

and it is this closure that makes the representation computationally
useful.

The consequences of this architectural choice are the subject of this
paper. Adding \(\mathrm{Si}(x)\) to the system required writing one
Taylor series generator and one line in the expansion engine --- not
because we found an implementation shortcut, but because
\(\mathrm{Si}(x)\) follows from the same representation as everything
else. The multivariate calculus requires seven lines of code, because
partial differentiation in the limit-first architecture is identical to
univariate differentiation treating other variables as constants ---
exactly the definition. The system returns \(\mathtt{NonElementary}\) as
a proved certificate when no elementary antiderivative exists, not as a
failure. And the product rule and chain rule are nowhere in the source
code, though both can be observed in the output, because they are
theorems about limits and our system computes limits.

The paper proceeds as follows. Section 2 defines the log-Puiseux
representation formally and establishes its closure properties. Section
3 develops the limit computation from this representation and states the
certified result type. Section 4 derives differentiation as coefficient
extraction and discusses the two-path architecture. Section 5 presents
the Risch integration algorithm as implemented over the log-Puiseux
representation, including the non-elementary certificate. Section 6
shows how multivariate calculus --- gradient, Jacobian, Hessian,
multivariate limits --- falls out of the architecture without additional
machinery. Section 7 states the open problems honestly, as a
mathematical research agenda. Section 8 concludes.




\section{The Log-Puiseux
Representation}\label{the-log-puiseux-representation}

We seek a class of formal series in \(h\) that is: (i) expressive enough
to represent the local behavior of all elementary functions and standard
special functions near any point; (ii) closed under the algebraic
operations needed to compute limits symbolically; and (iii) admits an
efficient finite representation.

The natural first candidate is the Puiseux series --- a generalization
of the Taylor series allowing rational exponents. We begin by showing
that this candidate is insufficient, and that the failure is not a
technicality but a structural fact about a specific and important class
of functions. This motivates the extension to log-Puiseux series, which
we then define and study.


\subsection{Why Pure Puiseux Series Are
Insufficient}\label{why-pure-puiseux-series-are-insufficient}

Recall that a \textbf{Puiseux series} in \(h\) centered at \(0\) is a
formal sum

\[\sum_{j=0}^{\infty} c_j \, h^{p_j}\]

where \(c_j \in \mathbb{C}\), \(p_j \in \mathbb{Q}\), and the exponents
\(p_j\) are bounded below and have a common denominator. These series
were introduced by Puiseux \cite{puiseux1850} in his study of algebraic
functions; the asymptotic ordering of such terms with respect to growth
rates is treated systematically in Hardy \cite{hardy1910}. A
\textbf{finite Puiseux series} is such a sum with only finitely many
nonzero terms. We work throughout with finite Puiseux series, which we
call \textbf{Puiseux polynomials} to emphasize finiteness.

The cosine integral is defined for \(x > 0\) by

\[\mathrm{Ci}(x) = -\int_x^\infty \frac{\cos t}{t}\,dt = \gamma + \log x + \int_0^x \frac{\cos t - 1}{t}\,dt\]

where \(\gamma\) is the Euler-Mascheroni constant. We wish to expand
\(\mathrm{Ci}(x + h)\) as a series in \(h\) near \(h = 0\).

\textbf{Lemma 2.1} \emph{(Necessity of logarithmic terms).} \emph{Let
\(x > 0\). The function \(h \mapsto \mathrm{Ci}(x+h)\) has no
representation as a Puiseux polynomial in \(h\) near \(h = 0\). More
precisely, for any finite collection of terms \(\sum_j c_j h^{p_j}\)
with \(c_j \in \mathbb{C}\) and \(p_j \in \mathbb{Q}\), the difference}

\[\mathrm{Ci}(x+h) - \sum_j c_j h^{p_j}\]

\emph{does not tend to zero as \(h \to 0^+\).}

\emph{Proof.} From the integral representation,

\[\mathrm{Ci}(x+h) = \gamma + \log(x+h) + \int_0^{x+h} \frac{\cos t - 1}{t}\,dt.\]

Expanding \(\log(x+h) = \log x + \log(1 + h/x)\) and writing the
integral as

\[\int_0^x \frac{\cos t - 1}{t}\,dt + \int_x^{x+h} \frac{\cos t - 1}{t}\,dt,\]

the second integral is \(O(h)\) as \(h \to 0\). Thus

\[\mathrm{Ci}(x+h) = \mathrm{Ci}(x) + \log\!\left(1 + \frac{h}{x}\right) + O(h).\]

Now \(\log(1 + h/x) = \frac{h}{x} - \frac{h^2}{2x^2} + \cdots\) is an
ordinary power series in \(h\), so this poses no difficulty for the
expansion at a fixed \(x > 0\). The issue arises at \(x = 0\). From the
second representation of \(\mathrm{Ci}\),

\[\mathrm{Ci}(h) = \gamma + \log h + \int_0^h \frac{\cos t - 1}{t}\,dt.\]

The integral
\(\int_0^h \frac{\cos t - 1}{t}\,dt = \int_0^h \frac{-t/2 + O(t^3)}{t}\,dt = -h/2 + O(h^3)\)
is a genuine power series in \(h\). But the term \(\log h\) is not.
Suppose for contradiction that \(\mathrm{Ci}(h) = \sum_j c_j h^{p_j}\)
for some finite Puiseux polynomial. Then

\[\log h = \mathrm{Ci}(h) - \gamma + h/2 - O(h^3) = \sum_j c_j h^{p_j} - \gamma + h/2 - O(h^3),\]

which would require \(\log h\) to be expressible as a finite sum of
terms \(h^{p_j}\) with rational exponents. But \(\log h \to -\infty\) as
\(h \to 0^+\) while every term \(h^{p_j}\) with \(p_j > 0\) tends to
zero and every term with \(p_j \leq 0\) either diverges as a power or is
constant. No finite combination of such terms can produce the
logarithmic divergence of \(\log h\), since
\(\log h = o(h^{-\epsilon})\) for every \(\epsilon > 0\) while
\(h^{p_j} \sim h^{p_j}\) exactly. The growth rates are incommensurable
in the sense of Hardy \cite{hardy1910}: \(\log h\) and any power
\(h^p\) with \(p < 0\) belong to different asymptotic scales, and no
finite linear combination of elements from one scale can reproduce the
behavior of the other.
\(\square\)

\textbf{Remark 2.2.} The same argument applies to
\(\mathrm{Ei}(h) = \gamma + \log h + h + O(h^2)\) and to
\(\mathrm{li}(h) = \mathrm{Ei}(\log h)\) near \(h = 1\). The logarithmic
singularity is not a pathology of a single function but a feature of the
class of functions definable by integrals whose integrands have
logarithmic antiderivatives --- precisely the functions the Risch
algorithm encounters.

Lemma 2.1 shows that pure Puiseux series cannot serve as our
representation class. The remedy is to add \(\log h\) as a primitive
alongside the power terms \(h^p\).


\subsection{Definition}\label{definition}

Let \(\bar{\mathbb{Q}}\) denote the algebraic closure of \(\mathbb{Q}\)
in \(\mathbb{C}\) --- the field of algebraic numbers. We represent
elements of \(\bar{\mathbb{Q}}\) as roots of minimal polynomials with
rational coefficients, identified by isolating rectangles in
\(\mathbb{C}\).

\textbf{Definition 2.3} \emph{(Log-Puiseux term).} A \textbf{log-Puiseux
term} in \(h\) is an expression of the form

\[c \cdot h^p \cdot \log(h)^k\]

where \(c \in \bar{\mathbb{Q}}\), \(p \in \mathbb{Q}\), and
\(k \in \mathbb{N}\).

\textbf{Definition 2.4} \emph{(Log-Puiseux series).} A
\textbf{log-Puiseux series} in \(h\) is a finite formal sum of
log-Puiseux terms:

\[S = \sum_{i=1}^{n} c_i \cdot h^{p_i} \cdot \log(h)^{k_i}\]

with \(c_i \in \bar{\mathbb{Q}} \setminus \{0\}\),
\(p_i \in \mathbb{Q}\), \(k_i \in \mathbb{N}\), and all pairs
\((p_i, k_i)\) distinct. The \textbf{zero series} is the empty sum. We
denote the set of all log-Puiseux series by \(\mathcal{L}\).

\textbf{Definition 2.5} \emph{(Support and leading term).} The
\textbf{support} of \(S \in \mathcal{L}\) is the set \(\{(p_i, k_i)\}\)
of exponent pairs with nonzero coefficients. We order pairs \((p, k)\)
lexicographically: \((p, k) < (p', k')\) if \(p < p'\), or \(p = p'\)
and \(k > k'\) (so higher log powers precede lower ones at the same
rational exponent, reflecting the fact that \(\log(h)^2 \to -\infty\)
faster than \(\log(h)\) as \(h \to 0^+\)). The \textbf{leading term} of
\(S\) is the term corresponding to the minimal element of the support
under this ordering.

\textbf{Remark 2.6.} The choice of coefficient field
\(\bar{\mathbb{Q}}\) rather than \(\mathbb{C}\) is deliberate. Algebraic
numbers admit exact arithmetic via resultant computations on minimal
polynomials; transcendental constants such as \(\pi\) or \(e\) cannot in
general be represented exactly in a finite symbolic system. The
constants that arise in calculus over elementary functions --- including
\(\gamma\), \(\sqrt{2}\), \(i\), roots of unity --- are algebraic or can
be carried symbolically as named constants. We return to this point when
discussing the implementation in Section 4.


\subsection{Algebraic Structure}\label{algebraic-structure}

We now establish that \(\mathcal{L}\) is closed under the operations we
need.

\textbf{Proposition 2.7} \emph{(Closure under addition).} \emph{If
\(S, T \in \mathcal{L}\), then \(S + T \in \mathcal{L}\).}

\emph{Proof.} Addition is termwise: collect terms sharing the same
exponent pair \((p, k)\) and add their \(\bar{\mathbb{Q}}\)
coefficients, discarding terms whose coefficients sum to zero. The
result is a finite formal sum of log-Puiseux terms with distinct
exponent pairs and nonzero coefficients, hence an element of
\(\mathcal{L}\). \(\square\)

\textbf{Proposition 2.8} \emph{(Closure under multiplication).} \emph{If
\(S, T \in \mathcal{L}\), then \(S \cdot T \in \mathcal{L}\).}

\emph{Proof.} Each term of \(S \cdot T\) is a product of one term from
\(S\) and one from \(T\):

\[(c_1 \cdot h^{p_1} \cdot \log(h)^{k_1}) \cdot (c_2 \cdot h^{p_2} \cdot \log(h)^{k_2}) = (c_1 c_2) \cdot h^{p_1 + p_2} \cdot \log(h)^{k_1 + k_2}.\]

Since \(\bar{\mathbb{Q}}\) is closed under multiplication,
\(c_1 c_2 \in \bar{\mathbb{Q}}\). Since \(\mathbb{Q}\) is closed under
addition, \(p_1 + p_2 \in \mathbb{Q}\). Since \(\mathbb{N}\) is closed
under addition, \(k_1 + k_2 \in \mathbb{N}\). The full product
\(S \cdot T\) is the sum over all such pairwise products; collecting
terms with equal exponent pairs and discarding zeros yields a finite sum
of log-Puiseux terms with distinct pairs. \(\square\)

\textbf{Corollary 2.9.} \((\mathcal{L}, +, \cdot)\) is a commutative
ring. It is not a field: the series \(h\) has no multiplicative inverse
in \(\mathcal{L}\) since \(h^{-1}\) is a valid log-Puiseux term but
\(h \cdot h^{-1} = 1\) requires the existence of a term with exponent
pair \((0, 0)\), which is present --- the issue is inversion of general
elements, not of monomials. More precisely, \(\mathcal{L}\) embeds into
the field of Puiseux series over \(\bar{\mathbb{Q}}(\log h)\), but
\(\mathcal{L}\) itself is a ring of truncated objects.

\textbf{Proposition 2.10} \emph{(Closure under composition, single
logarithm).} \emph{Let \(f\) be a function whose expansion \(f(x + h)\)
near \(h = 0\) lies in \(\mathcal{L}\), and let \(g\) be a function
whose expansion \(g(x + h)\) lies in \(\mathcal{L}\) with zero constant
term --- that is, \(g(x + h) = \sum_i c_i h^{p_i} \log(h)^{k_i}\) with
all \(p_i > 0\). Then the expansion of \(f(g(x+h))\) lies in
\(\mathcal{L}\).}

\emph{Proof sketch.} Write \(g(x+h) = u\) where \(u \to 0\) as
\(h \to 0\). Expand \(f(x + u)\) as a log-Puiseux series in \(u\), then
substitute \(u = \sum_i c_i h^{p_i} \log(h)^{k_i}\). Each power \(u^q\)
for \(q \in \mathbb{Q}\) becomes a finite product of log-Puiseux series
(by Proposition 2.8 applied finitely many times after truncation to a
fixed depth). The term \(\log(u)\) becomes

\[\log\!\left(\sum_i c_i h^{p_i} \log(h)^{k_i}\right) = \log(c_1 h^{p_1} \log(h)^{k_1}) + \log\!\left(1 + \frac{\text{lower order terms}}{c_1 h^{p_1} \log(h)^{k_1}}\right).\]

The first summand is \(\log c_1 + p_1 \log h + k_1 \log \log h\). This
last term, \(\log \log h\), lies outside \(\mathcal{L}\) in general.
However, when \(k_1 = 0\) --- that is, when the leading term of \(u\)
carries no logarithm --- we obtain
\(\log(c_1 h^{p_1}) = \log c_1 + p_1 \log h \in \mathcal{L}\). The
second summand's argument tends to \(0\) as \(h \to 0\), so it expands
as a convergent power series in the ratio, which after truncation lies
in \(\mathcal{L}\) by Proposition 2.8. Thus, provided the leading term
of \(g(x+h)\) is a pure power of \(h\) (no \(\log h\) factor), the
composition \(f(g(x+h))\) lies in \(\mathcal{L}\). \(\square\)

\textbf{Remark 2.11} \emph{(Boundary of the representation).}
Proposition 2.10 reveals an honest limitation: when the leading term of
the inner function itself carries a \(\log h\) factor, the composition
produces \(\log \log h\), which lies outside \(\mathcal{L}\). The
function \(\mathrm{li}(x) = \mathrm{Ei}(\log x)\) exhibits exactly this
behavior near \(x = 1\): expanding \(\mathrm{li}(1 + h)\) requires
\(\mathrm{Ei}(\log(1+h)) = \mathrm{Ei}(h - h^2/2 + \cdots)\), whose
expansion near \(h = 0\) involves \(\log \log h\). The system correctly
computes \(\frac{d}{dx}\mathrm{li}(x) = \frac{1}{\log x}\) via the
algebraic differentiation path (Section 4), but the series expansion
path returns \(\mathtt{LimitError}\) for \(\mathrm{li}\) at \(x = 0\).
This is a documented boundary of the representation, not an error.
Extending \(\mathcal{L}\) to include iterated logarithms is
mathematically straightforward in principle --- the type would carry a
finite tower of log levels --- but the composition algorithm for
\(\log \log h\) is not yet implemented.


\subsection{The Representation in
Practice}\label{the-representation-in-practice}

To make the definition concrete, we record the log-Puiseux expansions of
several functions that motivate and stress-test the representation.

\textbf{Example 2.12.} For \(x > 0\) fixed and \(h \to 0\):

\[\sin(x+h) = \sin x + (\cos x)\, h - \tfrac{1}{2}(\sin x)\, h^2 - \tfrac{1}{6}(\cos x)\, h^3 + \cdots\]

This is a pure power series (\(k_i = 0\) throughout); no log terms
appear. The representation reduces to an ordinary Taylor series with
algebraic-number coefficients at \(x\).

\textbf{Example 2.13.} At \(x = 0\):

\[\mathrm{Ci}(h) = \gamma + \log h - \tfrac{1}{4}h^2 + \tfrac{1}{96}h^4 - \cdots\]

The term \(\log h\) appears with coefficient \(1\) and exponent pair
\((0, 1)\); the constant \(\gamma\) has pair \((0, 0)\); the remaining
terms have pairs \((2,0), (4,0), \ldots\) This is genuinely a
log-Puiseux series and not expressible as a pure Puiseux series,
consistent with Lemma 2.1.

\textbf{Example 2.14} \emph{(Non-analyticity certificate).} The absolute
value function satisfies

\[|x + h| = |x| + \mathrm{sgn}(x) \cdot h \quad (x \neq 0)\]

which is a valid log-Puiseux expansion. At \(x = 0\), however, \(|h|\)
is not a log-Puiseux series in \(h\), since \(|h| = h\) for \(h > 0\)
and \(|h| = -h\) for \(h < 0\): the two-sided expansion does not exist
as a single series. The system returns \(\mathtt{NonAnalytic}\) in this
case --- a certified negative result.

\textbf{Example 2.15} \emph{(Fractional exponents).} The function
\(x^{1/2}\) at \(x = 1\) expands as

\[\sqrt{1+h} = 1 + \tfrac{1}{2}h - \tfrac{1}{8}h^2 + \tfrac{1}{16}h^3 - \cdots\]

with rational exponents \(p_i \in \{0, 1, 2, \ldots\}\) --- here
integer, as it happens --- and no log terms. At \(x = 0\), however,
\(\sqrt{h} = h^{1/2}\), a single Puiseux monomial with \(p = 1/2\),
\(k = 0\).




\section{Limits as the Primitive
Operation}\label{limits-as-the-primitive-operation}

The computational foundation of limcalc is limit evaluation. Every other
operation --- differentiation, integration --- is either derived from
limits directly or relies on the same series representation that makes
limit computation tractable. In this section we define the limit
computation formally, establish the certified result type, and prove the
properties that justify it.


\subsection{The Certified Result
Type}\label{the-certified-result-type}

A symbolic limit computation should not merely return a value when it
succeeds. It should communicate, in a machine-verifiable way, the nature
of the result: whether the limit exists, whether the function diverges
and at what rate, or whether the limit fails to exist for a stated
reason. We formalize this as a sum type.

\textbf{Definition 3.1} \emph{(Limit result).} The \textbf{limit result
type} \(\mathtt{LimitResult}\) is the disjoint union

\[\mathtt{LimitResult} = \mathtt{Exists}(\bar{\mathbb{Q}}) \mid \mathtt{Pole}(\mathbb{Q}^-) \mid \mathtt{DoesNotExist}(\mathtt{String}) \mid \mathtt{LimitError}(\mathtt{String})\]

where:

\begin{itemize}
\tightlist
\item
  \(\mathtt{Exists}(v)\) asserts that the limit exists and equals
  \(v \in \bar{\mathbb{Q}}\);
\item
  \(\mathtt{Pole}(r)\) asserts that \(f(x_0 + h)\) has leading
  log-Puiseux exponent \(r \in \mathbb{Q}^-\) as \(h \to 0^+\), meaning
  \(f\) diverges at rate \(h^r\) near \(x_0\);
\item
  \(\mathtt{DoesNotExist}(s)\) asserts that the limit does not exist,
  with \(s\) a human-readable witness;
\item
  \(\mathtt{LimitError}(s)\) indicates that the computation could not be
  completed within the current representation class.
\end{itemize}

\textbf{Remark 3.2} \emph{(On} \(\mathtt{Pole}(r)\)).* The rational
number \(r\) carried by \(\mathtt{Pole}\) is the leading exponent of the
log-Puiseux expansion of \(f(x_0 + h)\) --- formally, if the leading
term of the expansion is \(c \cdot h^r \cdot \log(h)^k\) then \(r\) is
that exponent. Concretely: \(1/x\) at \(x_0 = 0\) expands as \(h^{-1}\),
yielding \(\mathtt{Pole}(-1)\); \(1/x^2\) yields \(\mathtt{Pole}(-2)\);
\(1/\sqrt{x}\) yields \(\mathtt{Pole}(-1/2)\). The leading coefficient
\(c\) and log exponent \(k\) are not carried by the current
implementation --- the direction of divergence (whether
\(f \to +\infty\) or \(f \to -\infty\) as \(h \to 0^+\)) is not
recorded. This is a known gap: a complete treatment would carry the full
leading term \((c, r, k)\) rather than \(r\) alone. We note this
honestly and return to it in Section 7.

\textbf{Remark 3.3} \emph{(Semantics of} \(\mathtt{LimitError}\)).*
\(\mathtt{LimitError}\) is not a synonym for ``the limit does not
exist.'' It means the system was unable to determine the limit within
the log-Puiseux representation class --- for example, because the
function requires \(\log \log h\) terms (cf.~Remark 2.11). A
\(\mathtt{LimitError}\) result is epistemically neutral: the limit may
exist, may not exist, or may be a pole; the system simply cannot decide.
This distinction between certified non-existence
(\(\mathtt{DoesNotExist}\)) and computational incapacity
(\(\mathtt{LimitError}\)) is essential for the system to be
mathematically honest.




\subsection{Univariate Limit
Computation}\label{univariate-limit-computation}

Let \(f\) be a function whose local behavior near \(x_0\) is
representable in \(\mathcal{L}\). The computation proceeds by expanding
\(f(x_0 + h)\) as a log-Puiseux series and reading off the result from
the leading term.

\textbf{Definition 3.4} \emph{(Limit from the series).} Let
\(S = \sum_i c_i h^{p_i} \log(h)^{k_i} \in \mathcal{L}\) be the
log-Puiseux expansion of \(f(x_0 + h)\). Define
\(\mathtt{limitFromSeries}(S)\) as follows:

\begin{enumerate}
\def\labelenumi{\arabic{enumi}.}
\tightlist
\item
  If \(S\) is the zero series, return \(\mathtt{Exists}(0)\).
\item
  Let \((p^*, k^*)\) be the leading exponent pair of \(S\) (minimal
  under the ordering of Definition 2.5), with leading coefficient
  \(c^*\).
\item
  If \(p^* > 0\): the leading term vanishes as \(h \to 0\); return
  \(\mathtt{Exists}(c_0)\) where \(c_0\) is the coefficient of the
  \((0,0)\) term, or \(\mathtt{Exists}(0)\) if no such term is present.
\item
  If \(p^* = 0\) and \(k^* = 0\): the series has a nonzero constant
  term; return \(\mathtt{Exists}(c^*)\).
\item
  If \(p^* = 0\) and \(k^* > 0\): the leading term is
  \(c^* \log(h)^{k^*}\), which diverges logarithmically; return
  \(\mathtt{Pole}(0)\) --- a logarithmic divergence. (We extend
  \(\mathtt{Pole}\) to \(\mathbb{Q}^- \cup \{0\}\) for this case, since
  the divergence rate is \(\log(h)^{k^*}\) rather than a pure power.)
\item
  If \(p^* < 0\): return \(\mathtt{Pole}(p^*)\).
\end{enumerate}

\textbf{Theorem 3.5} \emph{(Correctness of univariate limit
computation).} \emph{Let \(f\) be a function analytic on a punctured
neighborhood of \(x_0\), and suppose \(f(x_0 + h)\) admits a log-Puiseux
expansion \(S \in \mathcal{L}\). Then \(\mathtt{limitFromSeries}(S)\)
correctly identifies the limit of \(f\) at \(x_0\) in the following
sense:}

\emph{(i) If \(\mathtt{limitFromSeries}(S) = \mathtt{Exists}(v)\), then
\(\lim_{h \to 0} f(x_0 + h) = v\).}

\emph{(ii) If \(\mathtt{limitFromSeries}(S) = \mathtt{Pole}(r)\) with
\(r < 0\), then \(|f(x_0 + h)| \sim |c^*| \cdot h^r\) as \(h \to 0^+\);
in particular \(|f(x_0 + h)| \to \infty\).}

\emph{Proof.} Both parts follow directly from the definition of the
log-Puiseux series as a formal expansion that converges in a punctured
neighborhood of \(h = 0\).

For (i): if the leading exponent \(p^* > 0\), then every term in \(S\)
vanishes as \(h \to 0\) except possibly the constant term (\(p = 0\),
\(k = 0\)). The value of the limit is therefore the constant coefficient
\(c_0\), or zero if none is present. If \(p^* = 0\) and \(k^* = 0\), the
leading term is the constant \(c^*\) and all remaining terms vanish,
giving limit \(c^*\).

For (ii): if \(p^* < 0\), the leading term \(c^* h^{p^*}\) dominates all
others as \(h \to 0^+\) (since \(p^* < p_i\) for all other terms \(i\),
so \(h^{p^*} \gg h^{p_i}\)). Thus \(f(x_0 + h) \sim c^* h^{p^*}\) and
\(|f(x_0 + h)| \to \infty\). \(\square\)

\textbf{Example 3.6.} We record several illustrative computations:

\begin{longtable}[]{@{}
  >{\raggedright\arraybackslash}p{(\columnwidth - 6\tabcolsep) * \real{0.2500}}
  >{\raggedright\arraybackslash}p{(\columnwidth - 6\tabcolsep) * \real{0.2500}}
  >{\raggedright\arraybackslash}p{(\columnwidth - 6\tabcolsep) * \real{0.2500}}
  >{\raggedright\arraybackslash}p{(\columnwidth - 6\tabcolsep) * \real{0.2500}}@{}}
\toprule\noalign{}
\begin{minipage}[b]{\linewidth}\raggedright
Function \(f\)
\end{minipage} & \begin{minipage}[b]{\linewidth}\raggedright
Point \(x_0\)
\end{minipage} & \begin{minipage}[b]{\linewidth}\raggedright
Expansion leading terms
\end{minipage} & \begin{minipage}[b]{\linewidth}\raggedright
Result
\end{minipage} \\
\midrule\noalign{}
\endhead
\bottomrule\noalign{}
\endlastfoot
\(\sin(x)/x\) & \(0\) & \(1 - h^2/6 + \cdots\) &
\(\mathtt{Exists}(1)\) \\
\(1/x\) & \(0\) & \(h^{-1}\) & \(\mathtt{Pole}(-1)\) \\
\(\mathrm{Ci}(x)\) & \(0\) & \(\gamma + \log(h) - h^2/4 + \cdots\) &
\(\mathtt{Pole}(0)\) (logarithmic) \\
\(x^{1/2}\) & \(0\) & \(h^{1/2}\) & \(\mathtt{Exists}(0)\) \\
\(|x|\) & \(0\) & not representable in \(\mathcal{L}\) &
\(\mathtt{LimitError}\) \\
\end{longtable}

Note that \(|x|\) at \(x_0 = 0\) in fact has a limit (\(0\)), but the
system cannot determine this from the series representation because
\(|h|\) is not a log-Puiseux series. This is an instance where
\(\mathtt{LimitError}\) is epistemically honest: the system declines to
answer rather than returning an incorrect result.




\subsection{Multivariate Limits via Path
Probing}\label{multivariate-limits-via-path-probing}

The univariate theory extends cleanly to the case where the limit point
is approached along a specified path. The multivariate limit problem ---
determining \(\lim_{(x,y) \to (x_0, y_0)} f(x,y)\) without specifying a
path --- is substantially harder. No complete symbolic decision
procedure is known for general multivariate limits of elementary
functions. The approach taken in limcalc is path-based: evaluate the
limit along a family of standard paths through the limit point and
synthesize the results.

\textbf{Definition 3.7} \emph{(Path family).} Fix a limit point
\((x_0, y_0) \in \mathbb{R}^2\). The \textbf{standard path family}
\(\mathcal{P}\) consists of the following seven parameterized curves,
each approaching \((x_0, y_0)\) as \(t \to 0^+\):

\begin{enumerate}
\def\labelenumi{\arabic{enumi}.}
\tightlist
\item
  Horizontal axis: \((x_0 + t,\ y_0)\)
\item
  Vertical axis: \((x_0,\ y_0 + t)\)
\item
  Diagonal \(y - y_0 = x - x_0\): \((x_0 + t,\ y_0 + t)\)
\item
  Anti-diagonal \(y - y_0 = -(x - x_0)\): \((x_0 + t,\ y_0 - t)\)
\item
  Line \(y - y_0 = m(x - x_0)\) for general slope \(m\):
  \((x_0 + t,\ y_0 + mt)\)
\item
  Parabola \(y - y_0 = (x - x_0)^2\): \((x_0 + t,\ y_0 + t^2)\)
\item
  Cubic \(y - y_0 = (x - x_0)^3\): \((x_0 + t,\ y_0 + t^3)\)
\end{enumerate}

Each path reduces the multivariate limit to a univariate limit in \(t\),
computed by the method of Section 3.2.

\textbf{Definition 3.8} \emph{(Path-based limit result).} Given
\(f(x,y)\) and limit point \((x_0, y_0)\):

\begin{enumerate}
\def\labelenumi{\arabic{enumi}.}
\tightlist
\item
  For each path \(\gamma \in \mathcal{P}\), compute
  \(L_\gamma = \mathtt{limitFromSeries}(f(\gamma(t)))\).
\item
  If any two paths return \(\mathtt{Exists}(v_1)\) and
  \(\mathtt{Exists}(v_2)\) with \(v_1 \neq v_2\), return
  \(\mathtt{DoesNotExist}\) with a witness identifying the two
  disagreeing paths.
\item
  If all paths return \(\mathtt{Exists}(v)\) for the same \(v\), return
  \(\mathtt{Exists}(v)\) as the candidate limit.
\item
  If any path returns \(\mathtt{Pole}\) or \(\mathtt{DoesNotExist}\),
  return accordingly.
\end{enumerate}

\textbf{Theorem 3.9} \emph{(Soundness of path disagreement).} \emph{If
two paths \(\gamma_1, \gamma_2 \in \mathcal{P}\) yield
\(\mathtt{Exists}(v_1)\) and \(\mathtt{Exists}(v_2)\) respectively with
\(v_1 \neq v_2\), then \(\lim_{(x,y) \to (x_0,y_0)} f(x,y)\) does not
exist.}

\emph{Proof.} Suppose for contradiction that
\(L = \lim_{(x,y) \to (x_0,y_0)} f(x,y)\) exists. Then for every
sequence \((x_n, y_n) \to (x_0, y_0)\), \(f(x_n, y_n) \to L\). In
particular, for any continuous path \(\gamma\) with
\(\gamma(0) = (x_0, y_0)\) and \(\gamma(t) \neq (x_0,y_0)\) for
\(t > 0\), the restriction \(f \circ \gamma\) satisfies
\(\lim_{t\to 0^+} f(\gamma(t)) = L\). Since both \(\gamma_1\) and
\(\gamma_2\) are such paths, we would have \(v_1 = L = v_2\),
contradicting \(v_1 \neq v_2\). \(\square\)

\textbf{Remark 3.10} \emph{(Incompleteness of path agreement).} Theorem
3.9 is a soundness result, not a completeness result. Agreement of all
seven paths on a common value \(v\) does not prove that
\(\lim_{(x,y)\to(x_0,y_0)} f(x,y) = v\). A function may return the same
limit along every path in \(\mathcal{P}\) while failing to have a limit
in the topological sense. The standard counterexample in two variables
is

\[f(x,y) = \frac{x^2 y}{x^4 + y^2}.\]

\textbf{Theorem 3.11} \emph{(The} \(x^2y/(x^4+y^2)\) \emph{pathology).}
\emph{The function \(f(x,y) = x^2y/(x^4+y^2)\) satisfies
\(\lim_{t\to 0} f(t, mt) = 0\) for every slope \(m \in \mathbb{R}\), and
\(\lim_{t\to 0} f(t, t^3) = 0\), yet \(\lim_{(x,y)\to(0,0)} f(x,y)\)
does not exist.}

\emph{Proof.} Along the line \(y = mx\):
\(f(t, mt) = t^2 \cdot mt / (t^4 + m^2t^2) = mt^3/(t^2(t^2 + m^2)) = mt/(t^2+m^2) \to 0\)
as \(t \to 0\).

Along the cubic \(y = t^3\):
\(f(t, t^3) = t^2 \cdot t^3/(t^4 + t^6) = t^5/(t^4(1+t^2)) = t/(1+t^2) \to 0\).

However, along the parabola \(y = x^2\):
\(f(t, t^2) = t^2 \cdot t^2/(t^4 + t^4) = t^4/(2t^4) = 1/2\).

Since the parabolic path yields limit \(1/2 \neq 0\), by Theorem 3.9 the
limit does not exist. \(\square\)

\textbf{Remark 3.12.} The system correctly handles this case: paths 1--5
(linear paths) and path 7 (cubic) return \(\mathtt{Exists}(0)\), while
path 6 (parabola \(y = x^2\)) returns \(\mathtt{Exists}(1/2)\). The
disagreement between paths triggers \(\mathtt{DoesNotExist}\) with the
parabolic path identified as the witness. This is precisely the kind of
certified negative result that distinguishes a limit-first architecture
from a system that merely evaluates along a default path.

\textbf{Remark 3.13} \emph{(Research agenda).} The completeness gap in
multivariate limit computation is an open mathematical problem, not an
implementation shortcoming. A complete decision procedure for
multivariate limits of elementary functions --- one that either proves
existence or certifies non-existence for any input --- is not known to
exist. The seven-path family is a principled choice: it covers all
linear approaches, a quadratic approach sufficient to detect the
classical counterexamples, and a cubic approach for higher-order
pathologies. Extending the path family to include algebraic curves of
arbitrary degree would improve coverage but would not achieve
completeness. We record the problem precisely: characterize the class of
functions \(f\) for which agreement on all polynomial paths of degree
\(\leq d\) implies existence of the limit. This appears to be an open
question in several complex variables.




\section{Differentiation as Coefficient
Extraction}\label{differentiation-as-coefficient-extraction}

The derivative of \(f\) at \(x\) is defined by

\[f'(x) = \lim_{h \to 0} \frac{f(x+h) - f(x)}{h}.\]

If \(f(x+h)\) admits a log-Puiseux expansion
\(S = \sum_i c_i h^{p_i} \log(h)^{k_i}\), then the constant term of
\(S\) is \(f(x)\), and the difference \(f(x+h) - f(x)\) is \(S\) with
the constant term removed. Dividing by \(h\) shifts every exponent by
\(-1\): the term \(c \cdot h^p \cdot \log(h)^k\) becomes
\(c \cdot h^{p-1} \cdot \log(h)^k\). The limit as \(h \to 0\) then
extracts the new constant term --- the coefficient of
\(h^1 \cdot \log(h)^0\) in the original expansion. That is:

\[f'(x) = \text{coefficient of } h^1 \log(h)^0 \text{ in the log-Puiseux expansion of } f(x+h).\]

This is the series path to differentiation. It requires no rules. It
requires no knowledge of the product rule or the chain rule. It requires
only that \(f(x+h)\) be expandable in \(\mathcal{L}\) and that we can
read off the \(h^1\) coefficient. In this section we make this precise,
establish its relationship to the standard algebraic derivative, and
show that the product rule and chain rule are consequences ---
observable in the output but nowhere declared in the system.




\subsection{The Algebraic Path: Differentiation in the Base
Field}\label{the-algebraic-path-differentiation-in-the-base-field}

We begin with the algebraic path, which is the primary implementation.

\textbf{Definition 4.1} \emph{(Derivation on expressions).} Let
\(\mathcal{E}\) denote the algebra of symbolic expressions over
\(\bar{\mathbb{Q}}(x)\) --- formal combinations of the variable \(x\),
algebraic constants, and the standard elementary functions and special
functions. Define
\(\mathtt{deriveExpr} : \mathcal{E} \times \mathtt{Var} \to \mathcal{E}\)
by structural recursion on the expression tree:

\[\mathtt{deriveExpr}(c, x) = 0 \quad (c \text{ constant})\]
\[\mathtt{deriveExpr}(x, x) = 1\]
\[\mathtt{deriveExpr}(x, y) = 0 \quad (y \neq x)\]
\[\mathtt{deriveExpr}(f + g,\ x) = \mathtt{deriveExpr}(f, x) + \mathtt{deriveExpr}(g, x)\]
\[\mathtt{deriveExpr}(f \cdot g,\ x) = \mathtt{deriveExpr}(f, x) \cdot g + f \cdot \mathtt{deriveExpr}(g, x)\]
\[\mathtt{deriveExpr}(f \circ g,\ x) = \mathtt{deriveExpr}(f, g) \cdot \mathtt{deriveExpr}(g, x)\]

together with the standard base cases for \(\sin\), \(\cos\), \(\exp\),
\(\log\), \(\mathrm{erf}\), \(\mathrm{Si}\), \(\mathrm{Ci}\),
\(\mathrm{Ei}\), \(\mathrm{arcsin}\), \(\mathrm{arctan}\), and so on.

\textbf{Definition 4.2} \emph{(Base field derivation).}
\(\mathtt{deriveBase} : \mathcal{E} \to \mathcal{E}\) is
\(\mathtt{deriveExpr}(\cdot, x)\) --- differentiation with respect to
the distinguished variable \(x\) in the base field
\(\bar{\mathbb{Q}}(x)\). For partial differentiation,
\(\mathtt{partialDiff}(f, y)\) is \(\mathtt{deriveExpr}(f, y)\),
treating all variables other than \(y\) as constants by construction.

\textbf{Remark 4.3.} The recursion clauses for \(f \cdot g\) and
\(f \circ g\) in Definition 4.1 are, respectively, the product rule and
the chain rule. They appear here as the definition of
\(\mathtt{deriveExpr}\), not as theorems about limits. This is the
standard algebraic approach: a derivation on an algebra is defined
axiomatically by the Leibniz rule \(D(fg) = D(f)g + fD(g)\), and the
chain rule is the corresponding statement for composition. The
relationship between this algebraic object and the analytic limit
definition is a theorem --- one that we will recover, rather than
assume, via the series path.




\subsection{The Series Path: Differentiation as Coefficient
Extraction}\label{the-series-path-differentiation-as-coefficient-extraction}

\textbf{Definition 4.4} \emph{(Symbolic log-Puiseux series).} A
\textbf{symbolic log-Puiseux series} is a log-Puiseux series whose
coefficients are elements of \(\mathcal{E}\) rather than
\(\bar{\mathbb{Q}}\) --- that is, expressions in \(x\) rather than
algebraic numbers. We denote the ring of symbolic log-Puiseux series by
\(\mathcal{L}_{\mathcal{E}}\).

\textbf{Definition 4.5} \emph{(Symbolic expansion).}
\(\mathtt{symExpand} : \mathcal{E} \to \mathcal{L}_{\mathcal{E}}\)
computes the symbolic log-Puiseux expansion of \(f(x+h)\) around \(x\),
to a fixed truncation depth \(N\). For the elementary functions, the
expansion is given by the Taylor series with symbolic coefficients:

\[\mathtt{symExpand}(\sin) = \sin(x) + \cos(x) \cdot h - \tfrac{\sin(x)}{2} \cdot h^2 - \tfrac{\cos(x)}{6} \cdot h^3 + \cdots\]
\[\mathtt{symExpand}(\exp) = e^x + e^x \cdot h + \tfrac{e^x}{2} \cdot h^2 + \cdots\]
\[\mathtt{symExpand}(\log) = \log(x) + \tfrac{1}{x} \cdot h - \tfrac{1}{2x^2} \cdot h^2 + \cdots\]

For the special functions with logarithmic singularities, the expansion
includes log terms as in Section 2.4.

\textbf{The key point} is how these coefficients are computed. The
\(n\)-th Taylor coefficient of \(f(x+h)\) around \(x\) is
\(f^{(n)}(x)/n!\), where \(f^{(n)}\) is the \(n\)-th derivative of
\(f\). In limcalc, this is computed by applying \(\mathtt{deriveBase}\)
iteratively:

\[[h^n] \,\mathtt{symExpand}(f) = \frac{1}{n!} \cdot \mathtt{deriveBase}^n(f)(x).\]

The series path is therefore not an independent differentiation
algorithm. It is the algebraic derivative, expressed through the Taylor
series representation. \(\mathtt{symExpand}\) does not know how to
differentiate \(\sin\) or \(\log\) or \(f \cdot g\) independently --- it
delegates entirely to \(\mathtt{deriveBase}\), which is
\(\mathtt{deriveExpr}\).

\textbf{Definition 4.6} \emph{(Differentiation via series extraction).}
\(\mathtt{diff}(f, x) = [h^1]\,\mathtt{symExpand}(f)\), where \([h^1]\)
denotes the coefficient of \(h^1 \log(h)^0\) in the symbolic expansion.




\subsection{The Two Paths Are One}\label{the-two-paths-are-one}

We now state precisely the relationship between the two paths.

\textbf{Theorem 4.7} \emph{(Series path as consequence of algebraic
path).} \emph{For any \(f \in \mathcal{E}\) whose Taylor series
converges in a neighborhood of \(x\),}

\[\mathtt{diff}(f, x) = \mathtt{deriveBase}(f)(x).\]

\emph{That is, the series path and the algebraic path compute the same
derivative.}

\emph{Proof.} By Definition 4.6,
\(\mathtt{diff}(f, x) = [h^1]\,\mathtt{symExpand}(f)\). By the
construction of \(\mathtt{symExpand}\) via iterated
\(\mathtt{deriveBase}\), the coefficient of \(h^1\) in the expansion is
\(\mathtt{deriveBase}^1(f)(x)/1! = \mathtt{deriveBase}(f)(x)\).
\(\square\)

Theorem 4.7 makes explicit what might otherwise seem like a coincidence:
the limit definition of the derivative and the algebraic derivation
operator give the same answer not because we arranged for them to agree,
but because the series expansion is implemented by iterating the
algebraic derivative. The two paths are not competing implementations of
the same operation. They are the same operation viewed from two levels
of abstraction: the series path presents the limit definition as a
computational procedure; the algebraic path presents the same
computation as structural recursion on the expression tree. The
implementation makes this identity explicit.

\textbf{Remark 4.8} \emph{(Why partialDiff uses the algebraic path
directly).} For iterated partial derivatives --- computing
\(\partial^2 f / \partial x \, \partial y\), for instance --- the series
path would attempt to expand an already-differentiated expression as a
log-Puiseux series in a new infinitesimal, introducing expansion
artifacts at symbolic points. The algebraic path is clean by
construction: \(\mathtt{partialDiff}(f, y)\) treats all variables other
than \(y\) as constants, because \(\mathtt{deriveExpr}(\cdot, y)\)
returns \(0\) for every variable except \(y\). Iterating is therefore:

\[\frac{\partial^2 f}{\partial x \, \partial y} = \mathtt{partialDiff}(\mathtt{partialDiff}(f, y),\ x)\]

which is two applications of \(\mathtt{deriveExpr}\) with different
variable arguments. No series expansion is involved, and no truncation
depth appears. The algebraic path is both simpler and more general for
the multivariate case.




\subsection{The Product Rule and Chain Rule as
Theorems}\label{the-product-rule-and-chain-rule-as-theorems}

We now demonstrate what the introduction promised: the product rule and
chain rule are not declared in the system but are observable as
consequences of the limit definition. The following two propositions are
theorems about log-Puiseux series, not axioms of the differentiation
algorithm.

\textbf{Proposition 4.9} \emph{(Product rule, recovered).} \emph{For
\(f, g \in \mathcal{E}\),}

\[\mathtt{diff}(f \cdot g,\ x) = \mathtt{diff}(f, x) \cdot g(x) + f(x) \cdot \mathtt{diff}(g, x).\]

\emph{Proof.} Expand \((f \cdot g)(x+h) = f(x+h) \cdot g(x+h)\) as a
product of symbolic log-Puiseux series. Writing
\(f(x+h) = f(x) + f'(x)h + O(h^2)\) and
\(g(x+h) = g(x) + g'(x)h + O(h^2)\), the product is

\[f(x)g(x) + (f'(x)g(x) + f(x)g'(x))h + O(h^2).\]

The coefficient of \(h^1\) is
\(f'(x)g(x) + f(x)g'(x) = \mathtt{diff}(f,x) \cdot g(x) + f(x) \cdot \mathtt{diff}(g,x)\).
\(\square\)

\textbf{Proposition 4.10} \emph{(Chain rule, recovered).} \emph{For
\(f, g \in \mathcal{E}\) with \(g\) differentiable at \(x\) and \(f\)
differentiable at \(g(x)\),}

\[\mathtt{diff}(f \circ g,\ x) = \mathtt{diff}(f,\ g(x)) \cdot \mathtt{diff}(g,\ x).\]

\emph{Proof.} Expand \(f(g(x+h))\) as a log-Puiseux series in \(h\).
Writing \(g(x+h) = g(x) + g'(x)h + O(h^2)\), set
\(u = g'(x)h + O(h^2)\), so \(g(x+h) = g(x) + u\) with \(u = O(h)\).
Then

\[f(g(x+h)) = f(g(x) + u) = f(g(x)) + f'(g(x)) \cdot u + O(u^2).\]

Substituting \(u = g'(x)h + O(h^2)\):

\[f(g(x+h)) = f(g(x)) + f'(g(x)) g'(x) h + O(h^2).\]

The coefficient of \(h^1\) is
\(f'(g(x)) \cdot g'(x) = \mathtt{diff}(f, g(x)) \cdot \mathtt{diff}(g, x)\).
\(\square\)

\textbf{Remark 4.11.} Propositions 4.9 and 4.10 are, of course, standard
results. What is nonstandard is where they appear in the logical
structure of the system. In a rule-table CAS, the product rule and chain
rule are axioms: the system applies them because it has been told to. In
limcalc, they are theorems: the system exhibits them because it computes
limits, and these are facts about limits. The source code contains
neither rule. Both can be observed in the output. This is the precise
sense in which the system is more mathematically honest --- it does not
invoke the product rule; it recovers it.




\subsection{Consistency and the Li
Asymmetry}\label{consistency-and-the-li-asymmetry}

Theorem 4.7 guarantees that the two paths agree whenever
\(\mathtt{symExpand}\) is defined. The known case where they diverge is
\(\mathrm{li}(x) = \mathrm{Ei}(\log x)\) near \(x = 0\).

\textbf{Proposition 4.12} \emph{(Derivative of} \(\mathrm{li}\)).*
\(\mathtt{deriveBase}(\mathrm{li}(x)) = 1/\log(x)\).

\emph{Proof.} By the chain rule applied algebraically:
\(\frac{d}{dx}\mathrm{Ei}(\log x) = \mathrm{Ei}'(\log x) \cdot \frac{1}{x} = \frac{e^{\log x}}{x} = \frac{x}{x \log x} \cdot \frac{1}{1}\).
More precisely, \(\mathrm{Ei}'(t) = e^t/t\), so
\(\frac{d}{dx}\mathrm{Ei}(\log x) = \frac{e^{\log x}}{\log x} \cdot \frac{1}{x} = \frac{x}{\log x \cdot x} = \frac{1}{\log x}\).
\(\square\)

However, \(\mathtt{diff}(\mathrm{li}(x), x)\) returns
\(\mathtt{LimitError}\), because
\(\mathtt{symExpand}(\mathrm{li}(x+h))\) near \(h = 0\) requires
expanding \(\mathrm{Ei}(\log(x+h))\), which near \(x = 0\) involves
\(\log \log h\) --- outside \(\mathcal{L}\) by Remark 2.11. The
algebraic path gives the correct answer; the series path correctly
reports that it cannot. This is not an inconsistency:
\(\mathtt{LimitError}\) does not assert a wrong value, it declines to
compute. The two paths are in agreement about what they know.

\textbf{Remark 4.13} \emph{(Reconciliation via}
\(\mathtt{deriveBase}\)).* For all functions other than \(\mathrm{li}\),
the series path was made consistent with the algebraic path by
generating the Taylor coefficients of \(\mathrm{erf}\), \(\mathrm{Si}\),
\(\mathrm{Ci}\), \(\mathrm{Ei}\) via iterated \(\mathtt{deriveBase}\)
rather than by independent Taylor series implementations. This
reconciliation --- replacing independent series implementations with
derived ones --- is architecturally significant: it ensures that the
system has one differentiation algorithm, not two, and that the series
expansion of any function is automatically consistent with its algebraic
derivative. The \(\mathrm{li}\) asymmetry is the unique remaining case
where the series representation is genuinely insufficient, and it is so
for a mathematically precise reason (Remark 2.11), not an implementation
gap.




\section{Integration: The Risch Algorithm and Non-Elementary
Certificates}\label{integration-the-risch-algorithm-and-non-elementary-certificates}

Integration enters the system differently from differentiation.
Differentiation, as Section 4 established, is coefficient extraction
from a log-Puiseux expansion --- a read operation on a representation
the system already computes. Integration is harder: given \(f\), find
\(F\) such that \(F' = f\). There is no analogous read operation. The
existence of an elementary antiderivative is not guaranteed, and
determining whether one exists is a nontrivial decision problem.

The answer to that decision problem is the Risch algorithm (Risch 1969),
which operates in the language of differential fields and decides, in
finite time, whether an elementary antiderivative exists and computes it
when it does. What is remarkable about the limit-first architecture is
that the log-Puiseux representation --- chosen for its role in
differentiation and limit computation --- is precisely the
representation in which the Risch algorithm is most naturally stated.
The differential field structure that Risch requires is already present.
Integration does not require a new layer; it inhabits the same
mathematical object.

This section also contains an honest admission. When the Risch algorithm
certifies that no elementary antiderivative exists, limcalc consults a
small table of known non-elementary antiderivatives --- the integrals
defining \(\mathrm{erf}\), \(\mathrm{Si}\), \(\mathrm{Ci}\),
\(\mathrm{Ei}\), and \(\mathrm{li}\). This table is a rules engine in
the precise sense the introduction criticized: these antiderivatives are
declared, not derived. We did teach the system that
\(\int e^{-x^2}\,dx = \frac{\sqrt{\pi}}{2}\mathrm{erf}(x)\), and
similarly for the other entries. The philosophical distinction we wish
to preserve is this: the \emph{derivatives} of these special functions
are not declared --- they are computed by \(\mathtt{deriveBase}\), which
derives them from the limit definition via the algebraic path. What is
declared is the inverse direction, the antiderivative. That inversion is
necessarily a declaration: no system can derive the name
``\(\mathrm{erf}\)'' from first principles; at some point you have to
say what you are calling things. The table is small, each entry is
verified by a differentiation check, and it is limited to functions the
log-Puiseux representation can already express. But it is a rules table,
and the paper says so.




\subsection{Differential Fields and the Setting for
Risch}\label{differential-fields-and-the-setting-for-risch}

The theory of differential fields and algebraic extensions was developed
by Kolchin \cite{kolchin1973}; the following definitions follow the
presentation in Bronstein \cite{bronstein2005}.

\textbf{Definition 5.1} \emph{(Differential field).} A
\textbf{differential field} \((K, D)\) is a field \(K\) equipped with a
map \(D : K \to K\) --- the \textbf{derivation} --- satisfying:

\[D(f + g) = D(f) + D(g) \quad \text{(additivity)}\]
\[D(f \cdot g) = D(f) \cdot g + f \cdot D(g) \quad \text{(Leibniz rule)}\]

for all \(f, g \in K\). The \textbf{constants} of \((K, D)\) are the
elements \(c \in K\) with \(D(c) = 0\).

\textbf{Example 5.2.} The field \(\bar{\mathbb{Q}}(x)\) of rational
functions over the algebraic numbers, with \(D = d/dx\), is a
differential field with constants \(\bar{\mathbb{Q}}\).

\textbf{Definition 5.3} \emph{(Elementary extension).} An
\textbf{elementary extension} of a differential field \((K, D)\) is a
differential field \((K(\theta), D)\) where \(\theta\) satisfies one of:

\begin{itemize}
\tightlist
\item
  \(D(\theta) = D(f)/f\) for some \(f \in K\) --- a \textbf{logarithmic
  extension}, \(\theta = \log f\);
\item
  \(D(\theta) = D(f) \cdot \theta\) for some \(f \in K\) --- an
  \textbf{exponential extension}, \(\theta = e^f\);
\item
  \(\theta\) is algebraic over \(K\).
\end{itemize}

The \textbf{elementary functions} over \(\bar{\mathbb{Q}}(x)\) are those
belonging to some tower of elementary extensions
\(\bar{\mathbb{Q}}(x) \subset K_1 \subset K_2 \subset \cdots \subset K_n\).

\textbf{Definition 5.4} \emph{(Elementary antiderivative).} \(f \in K\)
has an \textbf{elementary antiderivative} if there exists an elementary
extension \(K \subset L\) and \(F \in L\) with \(D(F) = f\).

The fundamental theorem underlying the Risch algorithm is Liouville's
structure theorem, which constrains the form any elementary
antiderivative must take.

\textbf{Theorem 5.5} \emph{(Liouville's theorem, (Liouville 1833; Risch
1969)).} \emph{Let \((K, D)\) be a differential field with algebraically
closed constant field \(C\), and let \(f \in K\). If \(f\) has an
elementary antiderivative, then there exist \(v \in K\),
\(c_1, \ldots, c_n \in C\), and \(u_1, \ldots, u_n \in K\) such that}

\[\int f = v + \sum_{i=1}^n c_i \log u_i.\]

\emph{That is, any elementary antiderivative of \(f\) over \(K\) lies in
\(K\) itself, extended by logarithms of elements of \(K\) --- no new
exponentials or algebraic elements are required.}

Liouville's theorem is the mathematical engine behind both the Risch
algorithm and the non-elementary certificate. It says: if an elementary
antiderivative exists, it has a very specific form. The Risch algorithm
searches for that form. If the search fails --- if no \(v\) and
\(c_i, u_i\) can be found satisfying the equation --- then no elementary
antiderivative exists, and this failure is a proof, not an accident.




\subsection{The Risch Algorithm in
limcalc}\label{the-risch-algorithm-in-limcalc}

The implementation follows Bronstein (Bronstein 2005) throughout. We
summarize the main cases and highlight the aspects specific to the
log-Puiseux setting.

\textbf{5.2.1 The Polynomial and Rational Cases}

For \(f \in \bar{\mathbb{Q}}[x]\) (polynomial), integration is termwise:
\(\int x^n\,dx = x^{n+1}/(n+1)\).

For \(f = p/q \in \bar{\mathbb{Q}}(x)\) (rational function), the
algorithm proceeds in two stages.

\textbf{Hermite reduction} splits \(f\) into a polynomial part, a proper
rational part with squarefree denominator, and a ``logarithmic part''
--- the piece whose antiderivative involves \(\log\) terms.
Specifically, if \(q = \prod_i q_i^{e_i}\) is the factorization of \(q\)
into squarefree factors, Hermite reduction produces
\(g \in \bar{\mathbb{Q}}(x)\) and \(h = r/s\) with \(s\) squarefree such
that

\[\int f = g + \int h.\]

The residual \(h\) has squarefree denominator and its antiderivative, if
elementary, must be a sum of logarithms.

\textbf{The Rothstein-Trager algorithm} then finds those logarithms. For
\(h = a/d\) with \(d\) squarefree, form the resultant polynomial

\[R(z) = \mathrm{res}_x(a - z \cdot d',\ d) \in \bar{\mathbb{Q}}[z].\]

\textbf{Theorem 5.6} \emph{(Rothstein-Trager, (Rothstein 1976; Trager
1976)).} \emph{The roots of \(R(z)\) are exactly the coefficients
\(c_i\) in the Liouville representation
\(\int h = \sum_i c_i \log u_i\), and the corresponding
\(u_i = \gcd(a - c_i d',\ d)\).}

\textbf{Corollary 5.7} \emph{(Non-elementary certificate for rational
functions).} \emph{\(h = a/d\) has no elementary antiderivative over
\(\bar{\mathbb{Q}}(x)\) if and only if some root of \(R(z)\) is not in
\(\bar{\mathbb{Q}}\).}

This corollary is the rational function case of the non-elementary
certificate. When \(R(z)\) has an irrational root, the system returns
\(\mathtt{NonElementary}\) --- a proved statement that no elementary
antiderivative exists, by Liouville's theorem and the completeness of
Rothstein-Trager.

\textbf{Example 5.8.} For \(f = 1/(x^2 + 1)\), the resultant
\(R(z) = \mathrm{res}_x(1 - z \cdot 2x,\ x^2+1)\) has roots
\(\pm i/2 \notin \bar{\mathbb{Q}} \cap \mathbb{R}\). Over
\(\mathbb{R}\), the system returns \(\mathtt{NonElementary}\). Over
\(\mathbb{C}\), the roots are algebraic and the antiderivative is
\(\frac{i}{2}\log(x-i) - \frac{i}{2}\log(x+i)\), which the system
returns before \(\mathtt{foldEuler}\) converts it to \(\arctan(x)\).




\subsection{Two Mathematical
Refinements}\label{two-mathematical-refinements}

Two edge cases in the Rothstein-Trager algorithm require mathematical
treatment beyond the standard presentation.

\textbf{5.3.1 The Logarithmic Derivative Degeneration}

\textbf{Proposition 5.9} \emph{(Logarithmic derivative pre-check).}
\emph{Let \(h = a/d\) with \(d\) squarefree. If \(a = c \cdot d'\) for
some \(c \in \bar{\mathbb{Q}}\), then \(\int h = c \log d\), and the
Rothstein-Trager resultant degenerates.}

\emph{Proof.} If \(a = c \cdot d'\), then
\(h = c \cdot d'/d = c \cdot D(\log d)\), so \(\int h = c \log d\)
directly. In the resultant computation, \(a - z \cdot d' = (c - z) d'\),
so \(\mathrm{res}_x((c-z)d', d)\). Since \(\gcd(d', d)\) may be
nontrivial (as \(d\) need not be squarefree in its derivative), the
resultant does not factor cleanly into the expected form --- the root
\(z = c\) causes \(\gcd(a - c \cdot d', d) = \gcd(0, d) = d\), returning
the entire denominator rather than its irreducible factors. This
produces incorrect output: for \(\int 2x/(x^2-1)\,dx\), the naive
Rothstein-Trager returns a sum of \(\log(1)\) terms, effectively zero,
rather than \(\log(x^2-1)\).

The remedy is a pre-check (\(\mathtt{logDerivativeCheck}\)) that detects
\(a = c \cdot d'\) before entering the resultant machinery and returns
\(c \log d\) directly. \(\square\)

\textbf{5.3.2 Numerical Root Rationalization}

The roots of the Rothstein-Trager resultant \(R(z)\) are algebraic
numbers, but the algorithm finds them numerically via Durand-Kerner
iteration. When \(R(z)\) has repeated rational roots, numerical
root-finding returns values close to but not equal to the true roots ---
for instance, \(0.5000000011\) and \(0.5000000003\) in place of two
copies of \(1/2\).

\textbf{Proposition 5.10} \emph{(Rationalization of Rothstein-Trager
roots).} \emph{When the input rational function has coefficients in
\(\mathbb{Q}\), the roots of the Rothstein-Trager resultant
\(R(z) \in \mathbb{Q}[z]\) that are rational are exactly representable
as fractions with bounded denominators. Specifically, if \(R(z)\) has a
rational root \(p/q\) in lowest terms, then \(q\) divides the leading
coefficient of \(R\).}

\emph{Proof.} By the rational root theorem, any rational root \(p/q\) in
lowest terms of a polynomial with integer coefficients satisfies
\(p \mid a_0\) and \(q \mid a_n\) where \(a_0, a_n\) are the constant
and leading coefficients. Since \(R(z) \in \mathbb{Q}[z]\), clearing
denominators gives an integer polynomial, and the bound on \(q\)
follows. \(\square\)

\textbf{Corollary 5.11.} After Durand-Kerner iteration, each numerical
root should be snapped to the nearest rational with denominator bounded
by a fixed constant (in the implementation, denominators \(\leq 100\)).
This rationalization step is not a heuristic --- it is justified by
Proposition 5.10, which guarantees that rational roots are exactly
representable as simple fractions when the input is rational.




\subsection{The Non-Elementary
Certificate}\label{the-non-elementary-certificate}

The Risch algorithm does not merely fail to find antiderivatives it
cannot compute --- it proves they do not exist. We state this precisely.

\textbf{Theorem 5.12} \emph{(Non-elementary certificate).} \emph{When
limcalc returns} \(\mathtt{NonElementary}\) \emph{for \(\int f\,dx\),
this is a proved mathematical statement: no elementary antiderivative of
\(f\) exists over \(\bar{\mathbb{Q}}(x)\) or its elementary extensions.}

\emph{Proof sketch.} The Risch algorithm is a decision procedure: for
\(f\) in the supported class of elementary functions, it either finds an
elementary antiderivative or proves none exists via the completeness of
Liouville's theorem and the Rothstein-Trager algorithm (rational case),
or the structure theorem for exponential extensions (exponential case).
The \(\mathtt{NonElementary}\) result is returned only when the decision
procedure has completed and determined non-existence. It is not returned
on timeout, truncation, or any other incomplete computation --- those
cases return \(\mathtt{LimitError}\) or are simply not reached.
\(\square\)

\textbf{Example 5.13.} The following integrals return
\(\mathtt{NonElementary}\), each with a proved certificate:

\begin{longtable}[]{@{}
  >{\raggedright\arraybackslash}p{(\columnwidth - 2\tabcolsep) * \real{0.5000}}
  >{\raggedright\arraybackslash}p{(\columnwidth - 2\tabcolsep) * \real{0.5000}}@{}}
\toprule\noalign{}
\begin{minipage}[b]{\linewidth}\raggedright
Integrand
\end{minipage} & \begin{minipage}[b]{\linewidth}\raggedright
Certificate
\end{minipage} \\
\midrule\noalign{}
\endhead
\bottomrule\noalign{}
\endlastfoot
\(e^{-x^2}\) & Hermite reduction produces no rational part; exponential
case Risch finds no solution in \(\bar{\mathbb{Q}}(x, e^{-x^2})\) \\
\(\sin(x)/x\) & Rational-in-exponential form via Euler; Risch finds no
elementary solution \\
\(\cos(x)/x\) & Same argument \\
\(e^x/x\) & Logarithmic derivative check passes; exponential Risch finds
no solution \\
\(1/\log(x)\) & Logarithmic case; Risch finds no solution in
\(\bar{\mathbb{Q}}(x, \log x)\) \\
\end{longtable}

These are not failures. They are mathematical results.




\subsection{Special Function Recognition: Integration Beyond the
Elementary}\label{special-function-recognition-integration-beyond-the-elementary}

When the Risch algorithm certifies non-existence of an elementary
antiderivative, a recognition layer checks whether the integrand matches
a known non-elementary pattern whose antiderivative is a named special
function. The definitions and properties of the special functions used
here --- $\mathrm{erf}$, $\mathrm{Si}$, $\mathrm{Ci}$, $\mathrm{Ei}$,
and $\mathrm{li}$ --- follow the standard references Abramowitz and
Stegun \cite{abramowitz1964} and Olver et al.\ \cite{olver2010}. We are
direct about what this layer is: a rules table. The antiderivatives of
these functions were encoded manually. The system was taught these
integrals in the same sense that a rule-table CAS is taught
\(\frac{d}{dx}\sin(x) = \cos(x)\).

What distinguishes this table from the rule tables the introduction
criticized is not its existence but its scope and its verification. Its
scope is small --- seven entries --- and bounded by the expressive
capacity of the log-Puiseux representation: only functions the system
can already represent as first-class objects with known series
expansions can appear as antiderivatives. Its verification is exact:
every entry is checked by applying \(\mathtt{deriveBase}\) to the
claimed antiderivative and confirming equality with the integrand. The
table does not assert correctness; it is proved correct by the algebraic
differentiation path.

\textbf{Definition 5.14} \emph{(Special function recognition).} After
the Risch algorithm determines that \(\int f\,dx\) has no elementary
antiderivative, the recognition layer checks \(f\) against the following
table of known non-elementary integrals:

\[\int e^{-x^2}\,dx = \frac{\sqrt{\pi}}{2}\,\mathrm{erf}(x)\]
\[\int \frac{1}{\log x}\,dx = \mathrm{li}(x)\]
\[\int \frac{\sin x}{x}\,dx = \mathrm{Si}(x)\]
\[\int \frac{\cos x}{x}\,dx = \mathrm{Ci}(x)\]
\[\int \frac{e^x}{x}\,dx = \mathrm{Ei}(x)\]
\[\int \frac{1}{\sqrt{1-x^2}}\,dx = \arcsin(x)\]
\[\int \frac{1}{1+x^2}\,dx = \arctan(x)\]

together with their compositions with linear functions \(ax + b\).

\textbf{Theorem 5.15} \emph{(Correctness of special function
recognition).} \emph{For each entry in Definition 5.14, the returned
special function \(F\) satisfies \(F' = f\) --- verified by applying
\(\mathtt{deriveBase}\) to \(F\) and confirming equality with \(f\).}

\emph{Proof.} Each identity is verified by differentiation, which is
exact in the algebraic path (Theorem 4.7). The recognition layer is not
a pattern-match that asserts correctness --- it is backed by a
differentiation check. \(\square\)

\textbf{Remark 5.16} \emph{(The asymmetry between differentiation and
integration).} The honest account of Section 5 is this: differentiation
of special functions is not declared --- it is derived.
\(\mathtt{deriveBase}(\mathrm{erf}(x)) = \frac{2}{\sqrt{\pi}}e^{-x^2}\)
is computed by the algebraic path, not looked up. Integration of special
functions is declared --- the antiderivatives are encoded as a table and
verified post-hoc by differentiation. This asymmetry is not a design
flaw; it reflects a genuine mathematical asymmetry. Differentiation is a
local operation fully determined by the limit definition, and the
limit-first architecture computes it without rules. Integration is a
global operation --- finding \(F\) such that \(F' = f\) requires knowing
not just the local behavior of \(f\) but what function has that
derivative everywhere --- and for non-elementary antiderivatives, no
finite derivation procedure can produce the name of that function. The
table is the irreducible cost of naming things.

What the limit-first architecture does provide is the right setting in
which to add new entries to the table cheaply: a function can be added
as a special function antiderivative if and only if it is already
representable in \(\mathcal{L}\), its derivative is computable by
\(\mathtt{deriveBase}\), and the differentiation check passes. In a
rule-table CAS, adding a new special function requires propagating new
rules through every module that must recognize it. Here, the cost is one
series generator, one table entry, and one check. The architecture did
not eliminate the rules table for integration; it minimized and bounded
it.

\textbf{Remark 5.17} \emph{(The Euler substitution and trigonometric
integration).} Trigonometric integrands are handled via the Euler
substitution \(e^{ix} = \cos x + i \sin x\), converting them to rational
functions in \(e^{ix}\) over \(\mathbb{C}\). The Risch exponential case
then applies. After integration, \(\mathtt{foldEuler}\) converts
Euler-form output back to \(\sin\) and \(\cos\). This post-processing
step is explicit and separate from the simplification loop --- wiring it
into general simplification would make every simplification trig-aware,
which is architecturally unclean. The separation is a deliberate design
choice: \(\mathtt{foldEuler}\) is called when and only when the result
of trigonometric integration needs to be presented in real form.




\subsection{The Partial Fraction Decomposition and Its Coupling
to
Risch}\label{the-partial-fraction-decomposition-and-its-coupling-to-risch}

A technical point worth recording precisely: partial fraction
decomposition and Risch integration are more tightly coupled than they
initially appear.

\textbf{Proposition 5.18} \emph{(Coupling of partial fractions and
Risch).} \emph{The same resultant and GCD machinery that computes the
Rothstein-Trager integration coefficients implicitly factors the
denominator polynomial. Partial fraction decomposition over
\(\bar{\mathbb{Q}}\) therefore requires polynomial factorization, which
in limcalc is accomplished as a consequence of running
Rothstein-Trager.}

The practical consequence: \(\mathtt{partialFractions}\) uses
Rothstein-Trager to identify irreducible factors, then applies the
residue formula

\[A_i = \frac{p(r_i)}{q(x)/(x - r_i)\big|_{x=r_i}}\]

for each linear factor at root \(r_i\) using the original numerator and
denominator. The factors come from the Rothstein-Trager run; the
coefficients come from the original rational function. These two
subproblems --- factoring the denominator and computing the numerator
coefficients --- must be kept logically separate: computing coefficients
from the sub-problems generated during factoring produces incorrect
results, as the residue formula applies to the original function, not
the sub-decompositions.




\section{Multivariate Calculus as a
Consequence}\label{multivariate-calculus-as-a-consequence}

The multivariate calculus --- partial derivatives, gradient, Jacobian,
Hessian --- is not a separate module in limcalc. It is not a
generalization requiring new machinery. It is a direct consequence of
the single operation \(\mathtt{partialDiff}(f, x_i)\) defined in Section
4, applied in different organizational patterns. This section makes that
precise and records the implementation sizes as mathematical evidence
for the claim.

Multivariate limits were developed in Section 3.3 via path-based
probing; we cross-reference that treatment rather than repeat it here.
The focus of this section is the differential calculus.




\subsection{Partial
Differentiation}\label{partial-differentiation}

\textbf{Definition 6.1} \emph{(Partial derivative).} For
\(f \in \mathcal{E}\) a symbolic expression in variables
\(x_1, \ldots, x_n\), the \textbf{partial derivative} with respect to
\(x_i\) is

\[\frac{\partial f}{\partial x_i} = \mathtt{partialDiff}(f, x_i) = \mathtt{deriveExpr}(f, x_i)\]

where \(\mathtt{deriveExpr}(\cdot, x_i)\) treats every variable \(x_j\)
with \(j \neq i\) as a constant, returning \(0\) for
\(\mathtt{deriveExpr}(x_j, x_i)\) when \(j \neq i\).

This is not a definition by analogy with the univariate case --- it is
the univariate case. \(\mathtt{deriveExpr}(f, x_i)\) does not know it is
computing a partial derivative. It is computing the derivative of \(f\)
with respect to \(x_i\), treating the other variables as elements of the
constant field. The partial derivative requires no new code, no new
concept, and no new module. It is the same operation as the univariate
derivative, applied to an expression that happens to contain other
variables.

\textbf{Proposition 6.2} \emph{(Correctness of partial
differentiation).} \emph{\(\mathtt{partialDiff}(f, x_i)\) computes
\(\partial f / \partial x_i\) in the analytic sense: for fixed \(x_j\)
\((j \neq i)\), \(\mathtt{partialDiff}(f, x_i)\) equals the derivative
of \(x_i \mapsto f(x_1, \ldots, x_n)\) as a univariate function.}

\emph{Proof.} By Theorem 4.7, \(\mathtt{deriveBase}\) computes the
correct univariate derivative for any \(f \in \mathcal{E}\).
\(\mathtt{partialDiff}(f, x_i)\) is \(\mathtt{deriveExpr}(f, x_i)\),
which by Definition 4.1 returns \(0\) for every occurrence of \(x_j\)
(\(j \neq i\)) and applies the standard differentiation rules to
\(x_i\). This is precisely the definition of the partial derivative:
differentiate with respect to \(x_i\) holding all other variables fixed.
\(\square\)

\textbf{Iterated partial derivatives} follow by applying
\(\mathtt{partialDiff}\) repeatedly:

\[\frac{\partial^2 f}{\partial x_j \, \partial x_i} = \mathtt{partialDiff}(\mathtt{partialDiff}(f, x_i),\ x_j).\]

As noted in Remark 4.8, the algebraic path handles iterated partials
cleanly because each application of \(\mathtt{deriveExpr}\) operates on
the expression tree directly, without series expansion. There are no
truncation artifacts, no depth parameters, and no special cases for
mixed partials.




\subsection{The Gradient, Jacobian, and
Hessian}\label{the-gradient-jacobian-and-hessian}

\textbf{Definition 6.3} \emph{(Gradient).} For
\(f : \mathbb{R}^n \to \mathbb{R}\), the \textbf{gradient} is

\[\nabla f = \left(\frac{\partial f}{\partial x_1},\ \ldots,\ \frac{\partial f}{\partial x_n}\right) = \mathtt{map}(\mathtt{partialDiff}(f,\ \cdot),\ [x_1, \ldots, x_n]).\]

The gradient is one line of code: map \(\mathtt{partialDiff}\) over the
list of variables.

\textbf{Definition 6.4} \emph{(Jacobian).} For
\(\mathbf{f} = (f_1, \ldots, f_m) : \mathbb{R}^n \to \mathbb{R}^m\), the
\textbf{Jacobian matrix} is

\[J_{\mathbf{f}} = \left[\frac{\partial f_i}{\partial x_j}\right]_{i,j} = \mathtt{map}(\mathtt{gradient}(\cdot),\ [f_1, \ldots, f_m]).\]

The Jacobian is one line of code: map the gradient over the list of
component functions.

\textbf{Definition 6.5} \emph{(Hessian).} For
\(f : \mathbb{R}^n \to \mathbb{R}\), the \textbf{Hessian matrix} is

\[H_f = \left[\frac{\partial^2 f}{\partial x_i \, \partial x_j}\right]_{i,j} = \mathtt{map}(\mathtt{gradient}(\cdot),\ \mathtt{gradient}(f)).\]

The Hessian is one line of code: apply the gradient to each component of
the gradient.

\textbf{Theorem 6.6} \emph{(The multivariate calculus in seven lines).}
\emph{The complete implementation of gradient, Jacobian, and Hessian in
limcalc --- including type signatures --- is seven lines of Haskell. We
reproduce it in full:}

\begin{Shaded}
\begin{Highlighting}[]
\OtherTok{gradient ::}\NormalTok{ [}\DataTypeTok{String}\NormalTok{] }\OtherTok{{-}\textgreater{}} \DataTypeTok{Expr} \OtherTok{{-}\textgreater{}}\NormalTok{ [}\DataTypeTok{Expr}\NormalTok{]}
\NormalTok{gradient vars f }\OtherTok{=} \FunctionTok{map}\NormalTok{ (\textbackslash{}v }\OtherTok{{-}\textgreater{}}\NormalTok{ partialDiff f v) vars}

\OtherTok{jacobian ::}\NormalTok{ [}\DataTypeTok{String}\NormalTok{] }\OtherTok{{-}\textgreater{}}\NormalTok{ [}\DataTypeTok{Expr}\NormalTok{] }\OtherTok{{-}\textgreater{}}\NormalTok{ [[}\DataTypeTok{Expr}\NormalTok{]]}
\NormalTok{jacobian vars fs }\OtherTok{=} \FunctionTok{map}\NormalTok{ (gradient vars) fs}

\OtherTok{hessian ::}\NormalTok{ [}\DataTypeTok{String}\NormalTok{] }\OtherTok{{-}\textgreater{}} \DataTypeTok{Expr} \OtherTok{{-}\textgreater{}}\NormalTok{ [[}\DataTypeTok{Expr}\NormalTok{]]}
\NormalTok{hessian vars f }\OtherTok{=}\NormalTok{ jacobian vars (gradient vars f)}
\end{Highlighting}
\end{Shaded}

\emph{These seven lines are not an optimized implementation. They are
the definitions. The gradient is the list of partial derivatives. The
Jacobian is the matrix of gradients of the component functions. The
Hessian is the Jacobian of the gradient. No additional machinery is
required because \(\mathtt{partialDiff}\) already does the work.}

We call this a theorem rather than a remark because it has mathematical
content: the fact that the multivariate calculus reduces to seven lines
is not a property of Haskell or of functional programming style. It is a
consequence of the architecture. In a rule-table system, multivariate
calculus requires independent machinery --- new rules for partial
derivatives, new data structures for matrix-valued outputs, new
simplification passes aware of the multivariate context. Here, none of
that is needed. The seven lines are all there is because partial
differentiation treating other variables as constants is not a special
case of the univariate derivative --- it is the univariate derivative,
applied to an expression that contains constants it hasn't been asked to
differentiate.




\subsection{Symmetry of Mixed
Partials}\label{symmetry-of-mixed-partials}

A natural question for any implementation of iterated partial
derivatives is whether it respects the symmetry of mixed partials. We
verify this directly.

\textbf{Proposition 6.7} \emph{(Schwarz's theorem, computational form).}
\emph{For \(f \in \mathcal{E}\) with continuous second partial
derivatives,}

\[\mathtt{partialDiff}(\mathtt{partialDiff}(f, x),\ y) = \mathtt{partialDiff}(\mathtt{partialDiff}(f, y),\ x)\]

\emph{as symbolic expressions, up to simplification.}

\emph{Proof.} Both sides equal
\(\partial^2 f / \partial x \, \partial y\) by Proposition 6.2 applied
twice. That the symbolic expressions are equal up to simplification
follows from the fact that \(\mathtt{deriveExpr}\) applies the same
algebraic rules in both orders, and the simplifier reduces both to the
same normal form. \(\square\)

\textbf{Example 6.8.} For \(f(x,y) = x^2 y + \sin(xy)\):

\[\frac{\partial f}{\partial x} = 2xy + y\cos(xy), \qquad \frac{\partial^2 f}{\partial y \, \partial x} = 2x + \cos(xy) - xy\sin(xy).\]

\[\frac{\partial f}{\partial y} = x^2 + x\cos(xy), \qquad \frac{\partial^2 f}{\partial x \, \partial y} = 2x + \cos(xy) - xy\sin(xy).\]

Both mixed partials agree, as Proposition 6.7 guarantees.




\subsection{A Demonstration: Non-Existence of a Multivariate
Limit}\label{a-demonstration-non-existence-of-a-multivariate-limit}

We close the section by connecting the multivariate differential
calculus to the multivariate limit theory of Section 3.3 through a
complete worked example that exercises both.

\textbf{Example 6.9.} Consider \(f(x,y) = x^2 y / (x^4 + y^2)\) from
Theorem 3.11. We compute the gradient at a point away from the origin,
then examine what happens as we approach the origin.

At a generic point \((x, y)\) with \(x^4 + y^2 \neq 0\):

\[\frac{\partial f}{\partial x} = \frac{2xy(x^4+y^2) - x^2y \cdot 4x^3}{(x^4+y^2)^2} = \frac{2xy^3 - 2x^5y}{(x^4+y^2)^2}\]

\[\frac{\partial f}{\partial y} = \frac{x^2(x^4+y^2) - x^2y \cdot 2y}{(x^4+y^2)^2} = \frac{x^6 - x^2y^2}{(x^4+y^2)^2}.\]

Both partial derivatives are well-defined away from the origin. At the
origin, however, the function itself does not have a limit (Theorem
3.11), so the gradient cannot exist there in any meaningful sense. The
system correctly computes the partial derivatives at generic points via
\(\mathtt{partialDiff}\) and correctly returns \(\mathtt{DoesNotExist}\)
for the limit at the origin via the path-probing algorithm of Section
3.3. The two computations are consistent: differentiability implies the
existence of the limit of the function, so non-existence of the limit
implies non-differentiability, which is reflected in the system's
behavior.

This example illustrates the coherence of the limit-first architecture
across its components: the limit machinery and the differential calculus
machinery are not independent modules that happen to agree --- they are
both consequences of the same underlying representation, and their
consistency is not arranged but unavoidable.




\section{Open Problems and Research
Agenda}\label{open-problems-and-research-agenda}

The limitations of limcalc are not uniformly the same kind of thing.
Some are implementation gaps --- places where more engineering would
close the distance between the current system and a complete realization
of the architecture. Others are mathematical obstacles --- places where
the architecture makes a problem precise that is not known to be
solvable. We distinguish these carefully, because the second kind are
more interesting than the first.




\subsection{The Log-Log Extension: Boundary of the
Representation}\label{the-log-log-extension-boundary-of-the-representation}

\textbf{The problem.} The log-Puiseux representation \(\mathcal{L}\)
cannot express \(\log \log h\). As established in Remark 2.11 and
Proposition 2.10, this means the series expansion path fails for
\(\mathrm{li}(x) = \mathrm{Ei}(\log x)\) near \(x = 0\), where the
expansion requires \(\log \log h\) terms. The algebraic path computes
\(\frac{d}{dx}\mathrm{li}(x) = 1/\log x\) correctly, but
\(\mathtt{diff}(\mathrm{li}(x), x)\) returns \(\mathtt{LimitError}\).

\textbf{The mathematical obstacle.} Extending \(\mathcal{L}\) to include
\(\log \log h\) is straightforward in principle --- the extended type
would carry terms \(c \cdot h^p \cdot \log(h)^k \cdot \log(\log(h))^j\)
with \(j \in \mathbb{N}\). The closure properties of Section 2.3
generalize: addition and multiplication remain closed, and the
composition theorem (Proposition 2.10) extends to the case where the
leading term of the inner function carries a \(\log h\) factor. The
obstacle is the composition algorithm: computing
\(\log(\log(h) + \text{lower order})\) requires expanding \(\log\)
around \(\log h\), which produces
\(\log \log h + \text{lower order} / \log h + \cdots\), a series in
\(1/\log h\) that does not terminate in any finite extension. A complete
treatment requires either truncating this expansion at a fixed depth ---
introducing an approximation --- or finding a closed-form representation
for the iterated logarithm tower that avoids the infinite regress.

\textbf{What a solution would require.} A composition algorithm for
\(\log \log h\) that is exact, terminating, and consistent with the
existing closure properties of \(\mathcal{L}\). More ambitiously: a
uniform treatment of iterated logarithm towers \(\log^{(k)} h\) as a
graded extension of \(\mathcal{L}\), with closure and composition
theorems at each level. This is a problem in the algebra of formal
series, not in the implementation.




\subsection{True Left-Hand Limits}\label{true-left-hand-limits}

\textbf{The problem.} The current implementation computes both
\(\lim_{x \to x_0^+}\) and \(\lim_{x \to x_0^-}\) by expanding from the
right --- substituting \(x_0 + h\) with \(h \to 0^+\) in both cases.
True left-hand limits require expanding \(f(x_0 - h)\) with
\(h \to 0^+\), which for functions with branch cuts or asymmetric
singularities produces different results.

\textbf{The consequence.} For functions where the left and right limits
differ --- \(\mathrm{sgn}(x)\) at \(0\), \(\lfloor x \rfloor\) at
integers, \(\log|x|\) at \(0\) --- the system returns the right-hand
limit in both cases rather than detecting the asymmetry and returning
\(\mathtt{DoesNotExist}\). This is a correctness gap in the certified
result type: \(\mathtt{DoesNotExist}\) should be returned when left and
right limits disagree, but the current implementation cannot detect the
disagreement.

\textbf{What a solution would require.} A parallel expansion engine that
substitutes \(x_0 - h\) and computes the log-Puiseux series in
\(h \to 0^+\). The algebraic machinery is unchanged; only the
substitution direction changes. The certified result type would then
compare left and right expansions and return \(\mathtt{DoesNotExist}\)
with a witness when they disagree. This is an implementation gap, not a
mathematical obstacle, and is the most straightforward of the open
problems to close.




\subsection{Adaptive Truncation
Depth}\label{adaptive-truncation-depth}

\textbf{The problem.} The log-Puiseux expansion engine uses a fixed
truncation depth \(N = 8\) --- it computes at most eight terms of the
expansion. For most functions this is sufficient, but for deeply
composed expressions --- \(\sin^{1/2}(\sin(x^2))\), for instance --- the
outer function may require more terms from the inner expansion than the
truncation provides, producing incorrect coefficients.

\textbf{The mathematical issue.} The appropriate depth is not fixed but
depends on the function: specifically, on how many leading terms cancel
in the expansion of \(f(x+h) - f(x)\) before the \(h^1\) coefficient
becomes visible. For functions with high-order contact at a point ---
functions that agree with their Taylor polynomial to high order --- the
depth must be correspondingly large to detect the leading term of the
difference.

\textbf{What a solution would require.} Adaptive depth: detect when
leading coefficients are zero and extend the expansion until a nonzero
coefficient is found or a provable bound on the vanishing order is
reached. The challenge is that ``leading coefficient is zero'' requires
exact zero detection in \(\bar{\mathbb{Q}}\), which is itself nontrivial
(cf.~the AlgNum zero detection problem discussed in the implementation).
A practical approach: expand to depth \(N\), check whether the relevant
coefficient is nonzero, and double the depth if not, up to a fixed
maximum. This is a heuristic with a known failure mode --- functions
that genuinely vanish to order greater than the maximum depth --- which
should be reported as \(\mathtt{LimitError}\) rather than silently
returning zero.




\subsection{Exact Algebraic Arithmetic for Transcendental
Constants}\label{exact-algebraic-arithmetic-for-transcendental-constants}

\textbf{The problem.} The current AlgNum implementation approximates
transcendental functions --- \(\sin\), \(\cos\), \(\exp\), \(\log\) ---
via Double floating-point arithmetic internally. As noted in the
implementation record, this means that
\(\mathtt{algExp}(\mathtt{algZero})\) returns approximately
\(1.000000238\) rather than exactly \(1\), due to accumulated rounding
in the Taylor series evaluation. The approximation is sufficient for the
series expansion engine in practice, but it introduces a principled gap
between the mathematical specification and the implementation.

\textbf{What a solution would require.} Genuine exact arithmetic for
transcendental constants would require one of two approaches. The first
is a full algebraic number tower in which \(\pi\), \(e\), and their
combinations are carried as symbolic constants with algebraic relations
--- a significant undertaking requiring a theory of transcendence degree
and algebraic independence that goes well beyond the current minimal
polynomial representation. The second, more tractable approach is to
keep polynomial and rational function coefficients as exact rationals
throughout and introduce floating-point approximation only at the final
evaluation step --- never during intermediate series computations. This
requires restructuring the coefficient arithmetic so that the Taylor
series for \(\sin(x+h)\) carries exact rational multiples of \(\sin(x)\)
and \(\cos(x)\) as symbolic constants, evaluated numerically only when a
specific point \(x\) is substituted. This is the architecturally correct
approach and corresponds to keeping the series symbolic rather than
numeric at the coefficient level.




\subsection{Multivariate Integration and the PDE
Obstruction}\label{multivariate-integration-and-the-pde-obstruction}

\textbf{The problem.} limcalc integrates univariate functions via the
Risch algorithm. It does not integrate multivariate functions. The
question the architecture makes precise is: for which
\(f(x_1, \ldots, x_n)\) does there exist \(F\) in the log-Puiseux
closure such that \(\partial F / \partial x_i = f\) for a specified
\(i\), and how do we find it?

\textbf{Why this is fundamentally harder than the univariate case.} The
univariate Risch algorithm succeeds because Liouville's theorem (Theorem
5.5) provides a complete structural characterization of elementary
antiderivatives: if one exists, it has a specific form, and the
algorithm searches that form exhaustively. No analogous structure
theorem is known for multivariate antiderivatives. The reason is that
multivariate integration is entangled with the theory of partial
differential equations in a way that univariate integration is not.

To see this concretely: finding \(F(x,y)\) such that
\(\partial F/\partial x = f(x,y)\) is a problem with a free function ---
integrating with respect to \(x\) leaves an arbitrary function of \(y\).
Finding \(F(x,y)\) such that \(\partial F/\partial x = f(x,y)\) and
\(\partial F/\partial y = g(x,y)\) simultaneously is an integrability
problem: it requires \(\partial f/\partial y = \partial g/\partial x\),
which is a PDE condition on the input. More generally, finding a
potential function for a vector field is the question of whether the
field is exact --- a question whose answer depends on the topology of
the domain as well as the algebraic structure of the functions.

At the level of generality needed for a symbolic decision procedure ---
arbitrary elementary functions of several variables --- multivariate
integration connects to central open problems in the symbolic solution
of PDEs. No complete algorithm is known. The Risch algorithm's success
in the univariate case depended on the fortunate coincidence that
differential fields, which are algebraic objects, capture exactly the
structure needed for the integration problem. In several variables, the
algebraic structure of differential rings does not alone determine
integrability; topological and analytic conditions intervene.

\textbf{What the architecture makes precise.} The log-Puiseux
representation makes the multivariate integration question precise in a
way that is new: we can ask, for functions in \(\mathcal{L}\), what is
the closure of \(\mathcal{L}\) under antidifferentiation in several
variables? This is a well-posed mathematical question. It does not have
a known answer. We record it as an open problem:

\textbf{Open Problem 7.1.} \emph{Characterize the class of functions
\(f \in \mathcal{L}\) in \(n \geq 2\) variables for which an
antiderivative \(F \in \mathcal{L}\) exists with
\(\partial F / \partial x_i = f\) for a given \(i\). Is this class
decidable? If so, give a decision procedure analogous to the Risch
algorithm.}

We do not know the answer. We believe the question is interesting.




\subsection{Completeness of Multivariate Limit Path
Probing}\label{completeness-of-multivariate-limit-path-probing}

This problem was stated precisely in Remark 3.13 and we record it here
for completeness.

\textbf{Open Problem 7.2.} \emph{Characterize the class of functions
\(f(x,y)\) for which agreement of \(\lim_{t\to 0} f(\gamma(t))\) across
all polynomial paths \(\gamma\) of degree \(\leq d\) implies existence
of \(\lim_{(x,y)\to(0,0)} f(x,y)\). Does such a characterization depend
on \(d\)? Is there a finite \(d\) for which the implication holds for
all elementary functions?}

The \(x^2y/(x^4+y^2)\) example (Theorem 3.11) shows that linear paths
(\(d=1\)) are insufficient. The parabolic path (\(d=2\)) detects that
example. Whether degree \(2\) is sufficient for all elementary
functions, or whether arbitrarily high degree paths are needed, is not
known. A positive answer for some finite \(d\) would give a complete
multivariate limit algorithm for the elementary functions; a negative
answer would establish that no finite path family suffices.




\subsection{Summary}\label{summary}

The open problems fall into two categories. Problems 7.2 (left-hand
limits) and 7.3 (adaptive depth) are implementation gaps: the
mathematical path to their resolution is clear and the obstacles are
engineering. Problems 7.1 (log-log extension) and 7.4 (exact AlgNum) are
mathematical gaps with known structure: the direction of the solution is
clear but the work is nontrivial. Problems 7.5 (multivariate
integration) and 7.6 (path probing completeness) are genuinely open
mathematical problems: the architecture makes them precise, but their
resolution requires mathematical progress that does not yet exist.

The third category is the most interesting. A system built on the right
abstraction does not merely solve the problems it was designed to solve
--- it makes visible the problems that remain. limcalc was designed to
compute calculus from first principles. What it revealed, in addition,
is a precise formulation of what multivariate symbolic integration would
require and a sharp statement of the completeness boundary for
multivariate limits. These are not failure modes. They are research
questions the architecture is responsible for having posed.




\section{Conclusion}\label{conclusion}

The rule-table approach to symbolic calculus was not a mistake. When the
major computer algebra systems were designed, computation was expensive
and the goal was correctness --- getting the right answer as efficiently
as possible. A table of differentiation rules, with the chain rule as a
meta-rule over compositions, achieves that goal elegantly. It is fast,
general, and correct. The systems built on that foundation are
extraordinary achievements, and nothing in this paper diminishes them.

What has changed is that computation is cheap. The cost that made the
rule-table approach necessary is no longer a binding constraint. And
when you revisit the foundational choice as a choice rather than a
necessity, a different question becomes available: not ``what is the
most efficient way to get the right answer,'' but ``what would a system
look like whose internal logic matches the mathematical reality it is
computing?'' The answer to that question is what this paper has tried to
build.

The distinction we have been drawing throughout is between correctness
and honesty. A system is correct if it returns the right answer. A
system is honest if it returns the right answer for the right reason ---
if the computation it performs corresponds to the mathematical
definition of the operation it claims to implement. Every mature CAS is
correct. limcalc is an attempt at honesty.

The derivative of \(f\) at \(x\) is a limit. limcalc computes it as a
limit. The product rule is a theorem about limits. limcalc treats it as
one --- it is nowhere in the source code, but it can be observed in the
output, because any system that computes limits correctly will exhibit
the product rule as a consequence. The chain rule is a theorem about
limits. limcalc treats it as one --- for the same reason. The
non-existence of an elementary antiderivative is a mathematical fact,
provable from Liouville's theorem. limcalc returns a certificate of that
fact, not a failure message. These are not implementation details. They
are consequences of having chosen the right foundation.

The foundation is the log-Puiseux series. This is the primary
mathematical contribution of the paper. The log-Puiseux type --- finite
formal sums of terms \(c \cdot h^p \cdot \log(h)^k\) with
\(c \in \bar{\mathbb{Q}}\), \(p \in \mathbb{Q}\), \(k \in \mathbb{N}\)
--- is the correct representation class for the local behavior of
elementary functions and standard special functions near any point. It
is closed under addition, multiplication, and composition (within the
single-logarithm tower). It is strictly more expressive than pure
Puiseux series, as the cosine integral demonstrates. And it is precisely
expressive enough: the boundary of the class --- where \(\log \log h\)
first appears --- corresponds exactly to the boundary of what the
limit-first architecture can compute, and that boundary is
mathematically sharp.

Everything else in the paper is a consequence of this choice.
Differentiation is coefficient extraction --- reading the \(h^1\) term
from an expansion the system already computes. The multivariate calculus
is seven lines of code, because partial differentiation treating other
variables as constants is not a special case of the univariate
derivative but the univariate derivative itself, applied to an
expression that contains constants it has not been asked to
differentiate. The Risch algorithm inhabits the same differential field
structure that the log-Puiseux representation naturally provides, so
integration --- including the non-elementary certificate --- requires no
new layer. The special function recognition table is the one honest
concession: we did teach the system that
\(\int e^{-x^2}\,dx = \frac{\sqrt{\pi}}{2}\mathrm{erf}(x)\), and we have
said so plainly. The derivatives of those special functions we did not
teach --- they are computed from the definition, the same as everything
else.

The open problems the architecture reveals are more interesting than the
problems it solves. Multivariate symbolic integration is not merely
unimplemented --- it is entangled with the symbolic solution of partial
differential equations in a way that has no known resolution. The
completeness of multivariate limit computation is not merely unproved
--- it connects to questions in several complex variables that appear to
be genuinely open. A system built on the right abstraction does not
merely solve the problems it was designed to solve. It makes visible the
problems that remain, and it makes them precise.

We close with the observation that motivated the work. A student
learning real analysis is told that the derivative is a limit, shown
that this implies the product rule and the chain rule, and then taught
to use those rules as the primary computational tools --- the limit
definition receding into the background, consulted for edge cases but
not for computation. This is pedagogically sensible. The rules are
faster to apply than the definition, and correctness is what the student
needs.

But something is lost in that recession. The rules, elevated to primary
status, begin to feel like the foundation rather than the consequence.
The limit definition begins to feel like a motivation rather than a
fact. A student who has only ever differentiated by rule application has
correct answers but may have a subtly wrong picture of what the
derivative is.

limcalc is, among other things, an argument that the picture does not
have to be wrong --- that a system can compute derivatives by actually
computing limits, that the rules can remain what they mathematically
are, which is theorems, and that the cost of this honesty, measured in
computation time and implementation complexity, is lower than one might
have expected.

The derivative is a limit. We computed it that way. The product rule and
chain rule appeared in the output. We never taught them to the system.




\begin{thebibliography}{99}

\bibitem{abramowitz1964}
Abramowitz, M., and Stegun, I.~A., eds. (1964).
\emph{Handbook of Mathematical Functions with Formulas, Graphs, and Mathematical Tables}.
National Bureau of Standards, Washington, D.C. Reprinted by Dover Publications, 1972.

\bibitem{bronstein2005}
Bronstein, M. (2005).
\emph{Symbolic Integration I: Transcendental Functions}, 2nd ed.
Springer-Verlag, Berlin.

\bibitem{davenport1988}
Davenport, J.~H., Siret, Y., and Tournier, \'{E}. (1988).
\emph{Computer Algebra: Systems and Algorithms for Algebraic Computation}.
Academic Press, London.

\bibitem{geddes1992}
Geddes, K.~O., Czapor, S.~R., and Labahn, G. (1992).
\emph{Algorithms for Computer Algebra}.
Kluwer Academic Publishers, Boston.

\bibitem{griewank2008}
Griewank, A., and Walther, A. (2008).
\emph{Evaluating Derivatives: Principles and Techniques of Algorithmic Differentiation}, 2nd ed.
SIAM, Philadelphia.

\bibitem{hardy1910}
Hardy, G.~H. (1910).
\emph{Orders of Infinity: The ``Infinit\"{a}rcalc\"{u}l'' of Paul du Bois-Reymond}.
Cambridge University Press, Cambridge.

\bibitem{kolchin1973}
Kolchin, E.~R. (1973).
\emph{Differential Algebra and Algebraic Groups}.
Academic Press, New York.

\bibitem{liouville1833}
Liouville, J. (1833).
M\'{e}moire sur la classification des transcendantes, et sur l'impossibilit\'{e} d'exprimer les racines de certaines \'{e}quations en fonction finie explicite des coefficients.
\emph{Journal de Math\'{e}matiques Pures et Appliqu\'{e}es}, 2, 56--105.

\bibitem{moses1971}
Moses, J. (1971).
Symbolic integration: The stormy decade.
\emph{Communications of the ACM}, 14(8), 548--560.

\bibitem{muller2001}
M\"{u}ller, N.~T. (2001).
The i{RRAM}: Exact arithmetic in {C}++.
In \emph{Computability and Complexity in Analysis},
Lecture Notes in Computer Science, vol.~2064, pp.~222--252.
Springer-Verlag, Berlin.

\bibitem{olver2010}
Olver, F.~W.~J., Lozier, D.~W., Boisvert, R.~F., and Clark, C.~W., eds. (2010).
\emph{{NIST} Handbook of Mathematical Functions}.
Cambridge University Press, Cambridge.

\bibitem{puiseux1850}
Puiseux, V. (1850).
Recherches sur les fonctions alg\'{e}briques.
\emph{Journal de Math\'{e}matiques Pures et Appliqu\'{e}es}, 15, 365--480.

\bibitem{risch1969}
Risch, R.~H. (1969).
The problem of integration in finite terms.
\emph{Transactions of the American Mathematical Society}, 139, 167--189.

\bibitem{rothstein1976}
Rothstein, M. (1976).
\emph{Aspects of Symbolic Integration and Simplification of Exponential and Primitive Functions}.
PhD thesis, University of Wisconsin-Madison.

\bibitem{trager1976}
Trager, B.~M. (1976).
Algebraic factoring and rational function integration.
In \emph{Proceedings of {SYMSAC} '76}, pp.~219--226. ACM Press.

\end{thebibliography}


\end{document}
